\documentclass[12pt,a4paper,UTF8]{ctexart}
\usepackage[a4paper,top=22mm,bottom=22mm,left=22mm,right=22mm,headheight=15pt]{geometry}
\usepackage{amsmath,amssymb,mathtools,bm}
\usepackage{graphicx}
\usepackage{xcolor}
\usepackage{booktabs,longtable,array}
\usepackage{enumitem}
\usepackage{fancyhdr}
\usepackage{caption}
\usepackage{float}
\usepackage{microtype}
\usepackage{setspace}
\usepackage[hidelinks,unicode]{hyperref}
\usepackage[most]{tcolorbox}

\definecolor{ContestBlue}{HTML}{17365D}
\definecolor{SoftBlue}{HTML}{EEF4FA}
\definecolor{SoftGray}{HTML}{F5F6F7}
\hypersetup{pdftitle={第 092 套},pdfauthor={}}

\setstretch{1.24}
\setlength{\parindent}{2em}
\setlength{\parskip}{0.35em}
\setlist{itemsep=0.35em,topsep=0.4em}
\captionsetup{font=small,labelfont=bf,justification=centering}
\setcounter{secnumdepth}{0}

\pagestyle{fancy}
\fancyhf{}
\fancyhead[L]{\small 第 092 套}
\fancyfoot[C]{\thepage}
\renewcommand{\headrulewidth}{0.4pt}

\ctexset{
  section={format=\Large\bfseries\color{ContestBlue},beforeskip=1.2em,afterskip=0.7em},
  subsection={format=\large\bfseries\color{ContestBlue},beforeskip=1.0em,afterskip=0.5em},
  subsubsection={format=\normalsize\bfseries,beforeskip=0.8em,afterskip=0.35em}
}

\makeatletter
\newcommand{\tightlist}{\setlength{\itemsep}{0pt}\setlength{\parskip}{0pt}}
\newcommand{\maxwidth}{\ifdim\Gin@nat@width>\linewidth\linewidth\else\Gin@nat@width\fi}
\newcommand{\maxheight}{\ifdim\Gin@nat@height>0.72\textheight 0.72\textheight\else\Gin@nat@height\fi}
\newcommand{\pandocbounded}[1]{\begingroup\setkeys{Gin}{width=\maxwidth,height=\maxheight,keepaspectratio}#1\endgroup}
\makeatother

\begin{document}

\begin{titlepage}
  \centering
  \vspace*{58mm}
  {\Huge\bfseries\color{ContestBlue} 第 092 套\par}
  \vfill
  \vspace*{28mm}
\end{titlepage}


\section{第 1 题}\label{ux7b2c-1-ux9898}

一个半径为 \(r\) 的自行车轮在水平地面上做纯滚动，角速度大小为
\(\omega\)。轮内均匀分布着 \(12\)
根轮辐，轮辐宽度可以忽略；除轮辐外，轮内为空。

\begin{enumerate}
\def\labelenumi{\arabic{enumi}.}
\item
  在高度

  \[y=r(1-\lambda),\qquad -1<\lambda<1\]

  处，沿垂直于轮面的方向射出一支箭，忽略重力造成的竖直位移。若入射点的横坐标可以合理选择，并要求无论箭速多小，箭都能避开轮辐穿过轮子，求
  \(\lambda\) 应满足的条件。
\item
  若 \(\lambda\) 不满足第 1 问的条件，例如 \(\lambda=0.3\)，则长度为
  \(l\) 的箭还必须具有足够大的速度。求箭的最小速度 \(v_m\)。
\end{enumerate}

\begin{figure}[H]
\centering
\includegraphics[width=0.858\linewidth,height=\textheight,keepaspectratio,alt={图 1-1　纯滚动车轮、固定射线高度及入射点的几何关系。}]{../figures/question-01.pdf}
\caption*{图 1-1　纯滚动车轮、固定射线高度及入射点的几何关系。}
\end{figure}

\subsection{一、题型定位}\label{ux4e00ux9898ux578bux5b9aux4f4d}

车轮的平移与转动同时发生，障碍物是整根轮辐。轮缘端点的摆线只能给出轮辐扫掠区的一部分；判断某一点会不会碰到轮辐，要把该高度与每根有限长轮辐的交点都算进去。有限长箭还要比较箭头至箭尾过面的时间与连续安全时间。

\subsection{二、知识点}\label{ux4e8cux77e5ux8bc6ux70b9}

\begin{itemize}
\tightlist
\item
  向右纯滚动时，轮心速度为 \(\omega r\)，轮体顺时针转动。
\item
  同一根轮辐与固定水平线的交点不固定在轮缘上。
\item
  相邻轮辐相位差 \(\Delta=\pi/6\)，不是 \(\pi/12\)。
\item
  有限长度条件给出交点参数的闭区间，区间内的驻点也须检查。
\item
  永久空白点取决于扫掠区间是否重叠；有限时间窗取决于按时间排序的全部碰撞事件。
\item
  箭长 \(l\) 对应过面时间 \(l/v\)，临界擦碰应排除。
\end{itemize}

\subsection{三、物理图像}\label{ux4e09ux7269ux7406ux56feux50cf}

取轮心在 \(t=0\) 时的地面投影为原点，设 \(\theta=\omega t\)。轮心为
\((r\theta,r)\)。先看轮心下方的水平线，\(0<\lambda<1\)。以一根初始朝下的轮辐为基准，用
\(u\) 表示它相对向下竖直线的顺时针转角。轮辐与水平线的交点相对轮心为
\((-\lambda r\tan u,-\lambda r)\)。

交点须在线段轮辐上。交点到轮心的距离不能超过 \(r\)，因此

\[\frac{\lambda r}{\cos u}\le r,\qquad |u|\le A=\arccos\lambda.\]

其地面横坐标为 \(r(u-\lambda\tan u)\)，再加上初相位对应的平移量。

图 1-2
上图画这一真实交点函数；两端表示交点在轮缘，内部极值表示轮辐中段扫到最外侧。下图把同一入射点的所有通过事件按时间排列，安全区间既可能夹在两次事件之间，也可能邻接车轮进入或离开该点的时刻。

\subsection{四、解题入口}\label{ux56dbux89e3ux9898ux5165ux53e3}

令

\[G_\lambda(u)=u-\lambda\tan u,\qquad |u|\le A.\]

先求 \(G_\lambda\) 的值域，再把它按相邻轮辐的平移间隔 \(\Delta\)
叠加。第二问固定入射点 \(p=x/r\)，逐支求

\[G_\lambda(u)=p-k\Delta,\qquad \theta=u+k\Delta,\]

其中 \(k\) 为整数。排序全部
\(\theta\)，取满足穿轮定义的最大无碰撞间隔。这样不会漏掉一根轮辐对同一点的多次通过。

\subsection{五、解析}\label{ux4e94ux89e3ux6790}

\begin{figure}[H]
\centering
\includegraphics[width=0.454\linewidth,height=\textheight,keepaspectratio,alt={图 1-2　整根轮辐的交点扫掠函数及固定入射点的完整碰撞时间序列。}]{../figures/solution-01.pdf}
\caption*{图 1-2　整根轮辐的交点扫掠函数及固定入射点的完整碰撞时间序列。}
\end{figure}

\textbf{（1）轮心下方的整根轮辐扫掠范围}

图 1-2 上图中横轴是轮辐角 \(u\)，纵轴是交点地面坐标
\(G_\lambda(u)\)。求导：

\[G_\lambda'(u)=1-\lambda\sec^2u.\]

内部驻点满足

\[\cos^2u=\lambda,\qquad u=\pm T,\qquad T=\arccos\sqrt\lambda.\]

因为 \(0<\lambda<1\)，\(T<A\)，这两个点确实在线段允许范围内。在 \(u=T\)
处，

\[G_\lambda(T)=T-\lambda\tan T
=\arccos\sqrt\lambda-\sqrt{\lambda(1-\lambda)}\equiv g(\lambda).\]

在 \(u=-T\) 处取相反值。轮缘端点给出的半宽为

\[G_\lambda(A)=\arccos\lambda-\sqrt{1-\lambda^2}.\]

它小于内部极值 \(g(\lambda)\)，因为 \(G\) 从 \(0\) 增至 \(T\) 后再降至
\(A\)。连续函数的整个值域因此为
\([-g,g]\)。每根轮辐扫过的是一个横坐标区间，而不只是两个端点轨迹上的点。

相邻轮辐的区间中心相差 \(\Delta=\pi/6\)。存在永久安全点的充要条件是

\[2g(\lambda)<\Delta.\]

临界方程为

\[\boxed{\arccos\sqrt{\lambda_*}-\sqrt{\lambda_*(1-\lambda_*)}=\frac\pi{12}}.\]

令 \(T=\arccos\sqrt{\lambda_*}\)，也可写成

\[T-\frac12\sin2T=\frac\pi{12},\qquad \lambda_*=\cos^2T.\]

\(g'(\lambda)=-\sqrt{(1-\lambda)/\lambda}<0\)，所以临界根唯一。数值为

\[T=0.76121465997,\qquad \boxed{\lambda_*=0.52417407552}.\]

若 \(\lambda>\lambda_*\)，可选

\[k\Delta+g(\lambda)<p<(k+1)\Delta-g(\lambda)\]

中的任一点。这个开区间在整个车轮经过期间都没有轮辐扫过。

仅用轮缘摆线会得到约
\(\lambda>0.3714\)；该范围内仍有轮辐中段进入所谓空白区，故它不是本题整根轮辐模型的正确条件。

\textbf{（2）轮心上方的安全范围}

令 \(\mu=-\lambda\in(0,1)\)。选初始朝上的轮辐，交点函数为

\[H_\mu(u)=u+\mu\tan u,\qquad |u|\le\arccos\mu.\]

\(H_\mu'(u)=1+\mu\sec^2u>0\)，所以最大扫掠半宽在轮缘端点，等于

\[\arccos\mu+\sqrt{1-\mu^2}.\]

永久安全条件是该半宽小于 \(\pi/12\)。令 \(a=\arccos\mu\)，边界满足

\[a+\sin a=\frac\pi{12}.\]

数值为

\[a=0.13108724816,\qquad \mu=0.99142036319.\]

因此第1问的完整答案为

\[\boxed{-1<\lambda<-0.99142036319
\quad\text{或}\quad 0.52417407552<\lambda<1}.\]

在临界等号处，相邻闭扫掠区间刚好接触；该点仍被轮辐碰到，不能算安全点。\(\lambda=0\)
时轮心本身经过每一个候选点，也没有永久空白点。

\textbf{（3）有限长箭的穿轮时间必须先定义}

第1问只按``避开轮辐''判断轮辐扫掠；若同时要求整支有限长箭始终处在轮内且避开轮缘，任意慢的穿越不可能，题面两问的约束须区分。第2问以下采用这一完整穿越约定，所得速度阈值是该约定下的结果。

具体把``整支箭穿过轮内开口''理解为：从箭头到箭尾经过轮面的全过程，该固定入射点都在圆轮内部，并且不遇到轮辐；轮缘本身也不能碰到。轮缘进入、离开该点的时刻限定允许区间。若允许箭的绝大部分在车轮尚未到来或已经离开以后通过，时间窗可任意延长，便不能由题面单独定义一个非零的最小速度。

对 \(\lambda=0.3\)，记

\[s=\sqrt{1-\lambda^2}=0.95393920142.\]

固定入射点 \(p=x/r\) 位于轮内的条件是

\[|p-\theta|<s,\qquad p-s<\theta<p+s.\]

由于轮辐初相位只影响 \(p\) 的平移，只需考察
\(0\le p<\Delta\)。每一次碰撞均满足

\[G(u)=p-k\Delta,\qquad \theta=u+k\Delta,\]

其中

\[G(u)=u-0.3\tan u,\qquad |u|\le A=\arccos0.3.\]

不能只保留一个求根支：\(G\) 在 \([-A,-T]\)、\([-T,T]\)、\([T,A]\)
上分别单调，其中

\[T=\arccos\sqrt{0.3}=0.99115658643.\]

先在三段内分别求所有根，再将碰撞时刻与轮缘进入、离开时刻合并。图 1-2
下图用这套完整序列表示安全窗。

\textbf{（4）比较全部时间窗}

该高度的内部极值和端点值为

\[g=0.53289901694,\qquad e=A-s=0.31216447136.\]

事件个数改变的入射相位由 \(\pm g\)、\(\pm e\) 对 \(\Delta\) 取余得到：

\[p=0.00930024134,\ 0.21143430424,\
0.31216447136,\ 0.51429853426.\]

在这些点之间，各事件连续且不会交换顺序。对任一支的根
\(u(p)\)，隐式微分给出

\[\frac{d\theta}{dp}=\frac{du}{dp}
=\frac1{1-0.3\sec^2u}.\]

故相邻事件间隔的驻点由两端导数相等求得；邻接轮缘的间隔则与
\(d(p\pm s)/dp=1\)
比较。逐段检查驻点和左右端极限，即可完成一维全局比较，而非任意选一根相邻轮辐。

最长安全窗邻接轮缘，其上确界出现在

\[p\longrightarrow p_*^+,\qquad p_*=g-\Delta=0.00930024134.\]

此时一对相切碰撞刚刚消失，剩余最后一次碰撞取中心单调支上的根

\[u_*-0.3\tan u_*=p_*,\qquad u_*=0.01328639414.\]

从这次碰撞后至轮缘离开的无量纲时间上确界为

\[\begin{aligned}
\alpha_{\rm sup}&=(p_*+s)-u_*\\
&=s-0.3\tan u_*\\
&=\boxed{0.94995304862}.
\end{aligned}\]

另一侧有镜像的等长窗口。只比较两次轮辐碰撞之间的内部窗口，得到上确界

\[\alpha_{\rm int}\simeq0.8048274.\]

它小于轮缘窗口。参考解只计算了一个局部分支，窗口约为 \(0.45427\)。

\textbf{（5）由时间窗得到速度条件}

箭头至箭尾过轮面的时间为 \(l/v\)，对应轮体转角
\(\omega l/v\)。排除碰到轮缘、轮辐的边界，要求

\[\frac{\omega l}{v}<\alpha_{\rm sup}.\]

因此严格安全速度满足

\[\boxed{v>1.05268360521\,\omega l}.\]

速度下确界为 \(v_{\rm inf}\simeq1.05268\,\omega l\)，等号不可取。

\newpage

\section{第 2 题}\label{ux7b2c-2-ux9898}

利用液体表面张力可以使一枚硬币浮在水面上。水的密度为
\(\rho_0\)，硬币的密度为 \(\rho\)，且 \(\rho>\rho_0\)；硬币半径为
\(r\)、厚度为 \(h\)，水的表面张力系数为
\(\sigma\)。在硬币边缘处，液面切线与水平面的夹角为 \(\theta_0\)。

\begin{enumerate}
\def\labelenumi{\arabic{enumi}.}
\tightlist
\item
  硬币上表面比远处水平液面低 \(l\)。只考虑硬币的受力平衡，用
  \(\theta_0\)、\(\sigma\)、\(g\)、\(\rho\)、\(\rho_0\)、\(r\)、\(h\)
  表示 \(l\)。
\item
  假设硬币较轻，使得 \(\theta_0\ll1\)。以竖直方向为 \(y\) 轴、水平方向为
  \(x\) 轴，取远处水面为 \(y=0\)。由液体的受力平衡写出液面截线 \(y(x)\)
  满足的微分方程。
\item
  联立前两问，求平衡液面 \(y(x)\) 以及 \(\theta_0\)，结果只含
  \(\rho\)、\(\rho_0\)、\(g\)、\(\sigma\)、\(r\)、\(h\)。
\end{enumerate}

进一步讨论：硬币较重时，\(\theta_0\)
最大只能达到固---液---气接触处的润湿角
\(\alpha\)。求硬币能够浮起的参数条件。

\begin{figure}[H]
\centering
\includegraphics[width=0.9\linewidth,height=\textheight,keepaspectratio,alt={图 2-1　硬币、下凹液面、浸没深度及边缘接触角。}]{../figures/question-02.pdf}
\caption*{图 2-1　硬币、下凹液面、浸没深度及边缘接触角。}
\end{figure}

\subsection{一、题型定位}\label{ux4e00ux9898ux578bux5b9aux4f4d-1}

硬币能否浮起先由整体竖直受力决定，液面下凹量随之确定。若要把这个结果推广到较大润湿角，还必须检查小坡度模型是否仍成立。

\subsection{二、知识点}\label{ux4e8cux77e5ux8bc6ux70b9-1}

\begin{itemize}
\item
  液体对水平底面的静水压力及阿基米德浮力的等价表达。
\item
  圆周接触线上的表面张力竖直分量 \(2\pi r\sigma\sin\theta_0\)。
\item
  杨---拉普拉斯关系与小坡度线性化。
\item
  毛细长度

  \[a=\sqrt{\frac{\sigma}{\rho_0g}}.\]
\item
  二阶常系数微分方程的指数衰减解。
\item
  边界条件 \(y(\infty)=0\)、\(y(r)=-l\) 和 \(y'(r)\simeq\theta_0\)。
\item
  润湿角给出的最大表面张力支撑条件。
\end{itemize}

\subsection{三、物理图像}\label{ux4e09ux7269ux7406ux56feux50cf-1}

硬币比水密度大，但不意味着只能下沉。硬币受到三类竖直作用：重力向下；水对硬币底面的静水压力向上；接触线处表面张力沿液面切线作用，其竖直分量也向上。

远处液面取 \(y=0\)，硬币上表面位于 \(y=-l\)，底面位于
\(y=-(l+h)\)。液面从硬币边缘的 \(-l\)
逐渐恢复到零。重力倾向于把被压低的液面拉回水平，表面张力则抑制曲率；两者竞争产生特征长度
\(a\)。

\subsection{四、解题入口}\label{ux56dbux89e3ux9898ux5165ux53e3-1}

第一问必须先画出硬币整体受力图，特别注意静水压力对应的是硬币底面深度
\(l+h\)，而不是 \(l\)。第二问先在二维小坡度近似下得到

\[\sigma y''=\rho_0gy,\]

再由远处边界条件选择衰减指数。第三问用

\[y(r)=-l,
\qquad
\theta_0\simeq y'(r)
\]

把未知的 \(l\) 与 \(\theta_0\) 消去。

\subsection{五、解析}\label{ux4e94ux89e3ux6790-1}

图 2-2
左侧是硬币整体的竖直受力图：重力、浮力和圆周表面张力的竖直分量同时出现；右侧是从硬币边缘
\(x=r\) 向远处恢复的液面曲线，并标出毛细长度
\(a=\sqrt{\sigma/(\rho_0g)}\)。第（1）步使用左图，第（2）---（4）步使用右图。

\begin{figure}[H]
\centering
\includegraphics[width=0.9\linewidth,height=\textheight,keepaspectratio,alt={图 2-2　硬币整体受力与小坡度毛细面的指数恢复。}]{../figures/solution-02.pdf}
\caption*{图 2-2　硬币整体受力与小坡度毛细面的指数恢复。}
\end{figure}

\textbf{（1）硬币整体的竖直受力平衡} 按图 2-2
左侧，向下只有硬币重力，向上有浮力和沿圆周分布的表面张力竖直分量。逐项写出后得到

硬币体积为

\[V=\pi r^2h,\]

所以重力为

\[G=\rho\pi r^2hg.\]

以大气压为零表压。硬币底面位于远处自由液面以下 \(l+h\)，底面处水的表压为

\[p_b=\rho_0g(l+h).\]

由于底面水平，水对底面的向上压力为

\[F_p=p_b\pi r^2=\rho_0g(l+h)\pi r^2.\]

硬币上表面接触空气，大气压在上下表面的共同部分已经相消，故不再另列一项大气压力。

接触线是半径为 \(r\) 的圆，其周长为

\[L_c=2\pi r.\]

单位长度接触线上的表面张力大小为 \(\sigma\)，沿液面切线方向。竖直分量为
\(\sigma\sin\theta_0\)，故总竖直支撑力为

\[F_\sigma=2\pi r\sigma\sin\theta_0.\]

平衡时向上的力等于重力：

\[
\rho_0g(l+h)\pi r^2
+2\pi r\sigma\sin\theta_0
=\rho\pi r^2hg.
\]

两边除以 \(\pi r^2\rho_0g\)：

\[
l+h+\frac{2\sigma\sin\theta_0}{r\rho_0g}
=\frac{\rho}{\rho_0}h.
\]

把 \(l\) 单独留在左边：

\[
\begin{aligned}
l
&=\frac{\rho}{\rho_0}h-h
-\frac{2\sigma\sin\theta_0}{r\rho_0g}\\
&=\frac{\rho-\rho_0}{\rho_0}h
-\frac{2\sigma\sin\theta_0}{r\rho_0g}.
\end{aligned}
\]

因此

\[
\boxed{
l=\frac{(\rho-\rho_0)h}{\rho_0}
-\frac{2\sigma\sin\theta_0}{r\rho_0g}
}.
\]

检查量纲：第一项是长度；第二项中 \(\sigma/(\rho_0gr)\)
的量纲也是长度，结果自洽。

\textbf{（2）建立液面微分方程} 转看图 2-2 右侧。取 \(x\)
轴沿未受扰动液面、\(y\) 轴向上，液面低于远处水平面时
\(y<0\)；曲线向右逐渐恢复到 \(y=0\)。

先采用参考解答中的局部二维模型，即再假设
\(r\gg a\)，忽略绕硬币圆周的另一主曲率。考虑液面上一小段弧，其单位垂直纸面宽度内的两端表面张力合力，法向分量等于
\(\sigma\kappa\,\mathrm{d}s\)，其中平面曲线曲率为

\[
\kappa=\frac{y''}{(1+y'^2)^{3/2}}.
\]

该处液面相对远处下降到 \(y<0\)，静水压相对于远处的增量大小为
\(-\rho_0gy\)。选取符号使下凹的 \(y\) 为负，法向平衡可写成

\[
\sigma\frac{y''}{(1+y'^2)^{3/2}}=\rho_0gy.
\]

由于 \(\theta_0\ll1\)，整段液面的坡度很小：

\[|y'|\ll1.\]

于是

\[
(1+y'^2)^{3/2}\simeq1,
\]

微分方程线性化为

\[
\boxed{\sigma y''=\rho_0gy}.
\]

也可写成

\[
\boxed{y''-\frac{1}{a^2}y=0},
\qquad
a=\sqrt{\frac{\sigma}{\rho_0g}}.
\]

毛细长度 \(a\) 表示液面从边缘扰动恢复到水平的大致距离。

\textbf{（3）求液面函数 \(y(x)\)} 右图表明物理解必须随 \(x-r\)
增大而衰减，因此通解中的增长指数项必须舍去。

常系数方程的特征方程为

\[k^2-\frac1{a^2}=0,\]

所以

\[k=\pm\frac1a.\]

通解为

\[
y(x)=A\exp\left(\frac{x-r}{a}\right)
+B\exp\left(-\frac{x-r}{a}\right).
\]

远离硬币时液面回到 \(y=0\)：

\[\lim_{x\to\infty}y(x)=0.\]

第一项随 \(x\) 指数增大。要满足远处有限且趋零，必须有

\[A=0.\]

于是

\[y(x)=B\exp\left(-\frac{x-r}{a}\right).\]

在硬币边缘 \(x=r\) 处，液面与硬币上表面等高：

\[y(r)=-l.\]

代入得到

\[B=-l.\]

因此

\[
\boxed{
y(x)=-l\exp\left[-\frac{x-r}{a}\right]
},
\qquad x\ge r.
\]

求导：

\[
y'(x)=\frac{l}{a}
\exp\left[-\frac{x-r}{a}\right].
\]

在边缘处

\[y'(r)=\frac la.\]

小角条件下 \(\tan\theta_0\simeq\theta_0\)，故

\[
\boxed{\theta_0\simeq\frac la}.
\]

\textbf{（4）联立硬币平衡条件，消去 \(l\) 与 \(\theta_0\)}

在小角近似下

\[\sin\theta_0\simeq\theta_0\simeq\frac la.\]

代回第一问结果：

\[
l=\frac{(\rho-\rho_0)h}{\rho_0}
-\frac{2\sigma}{r\rho_0g}\frac la.
\]

利用

\[a^2=\frac{\sigma}{\rho_0g},\]

第二项化为

\[
\frac{2\sigma}{r\rho_0g}\frac la
=\frac{2a^2}{r}\frac la
=\frac{2a}{r}l.
\]

所以

\[
l=\frac{(\rho-\rho_0)h}{\rho_0}-\frac{2a}{r}l.
\]

把含 \(l\) 的项移到左边：

\[
l\left(1+\frac{2a}{r}\right)
=\frac{(\rho-\rho_0)h}{\rho_0}.
\]

解得

\[
\boxed{
l=
\frac{(\rho-\rho_0)h}
{\rho_0\left(1+2a/r\right)}
}.
\]

于是边缘角为

\[
\boxed{
\theta_0=\frac la
=\frac{(\rho-\rho_0)h}
{\rho_0a\left(1+2a/r\right)}
}.
\]

最终液面为

\[
\boxed{
y(x)=
-\frac{(\rho-\rho_0)h}
{\rho_0\left(1+2a/r\right)}
\exp\left[-\frac{x-r}{a}\right]
},
\]

其中

\[
\boxed{a=\sqrt{\frac{\sigma}{\rho_0g}}}.
\]

\textbf{（5）润湿角限制与小坡度模型的适用范围}

真实圆硬币的液面为轴对称曲面。小坡度的严格径向方程为

\[y''+\frac{y'}x-\frac y{a^2}=0,\]

远处衰减解与边缘斜率为

\[y(x)=-l\frac{K_0(x/a)}{K_0(r/a)},\qquad
\theta_0\simeq\frac la\frac{K_1(r/a)}{K_0(r/a)}.\]

令 \(D=(\rho-\rho_0)h/\rho_0\)，与硬币整体平衡联立得到

\[\boxed{l=\frac D{1+(2a/r)K_1(r/a)/K_0(r/a)}}.\]

这是小角但不要求大半径的结果。前面的指数式还需要 \(r\gg a\)，使方位曲率
\(y'/x\) 可略；不能仅由 \(\theta_0\ll1\) 就删去它。

硬币较重时不能把 \(l=a\theta_0\) 继续用于任意润湿角。若采用 \(r\gg a\)
的二维近似，应保留完整曲率方程

\[\frac{y''}{(1+y'^2)^{3/2}}=\frac y{a^2}.\]

两边乘 \(y'\) 并从远处水平液面积分，得到

\[\frac1{\sqrt{1+y'^2}}=1-\frac{y^2}{2a^2}.\]

在边缘 \(y=-l\)、\(y'=\tan\theta_0\)，且 \(0\le\theta_0\le\pi/2\)，故

\[\boxed{l=2a\sin(\theta_0/2)}.\]

代回整体平衡，重硬币的二维近似条件为

\[D=2a\sin\frac{\theta_0}2+\frac{2a^2}r\sin\theta_0.\]

右端在 \(0\le\theta_0\le\alpha\le\pi/2\) 单调增加。因此能够支撑的条件为

\[\boxed{D\le2a\sin\frac\alpha2+\frac{2a^2}r\sin\alpha},\qquad r\gg a.\]

\(\alpha\ll1\) 时才退化为 \(D\le a\alpha(1+2a/r)\)。若 \(r\) 与 \(a\)
可比且角度较大，须求解轴对称完整曲率方程

\[\frac{y''}{(1+y'^2)^{3/2}}+
\frac{y'}{x\sqrt{1+y'^2}}=\frac y{a^2}\]

与
\(y(\infty)=0\)、\(y'(r)=\tan\theta_0\)，再代入硬币平衡式求最大载重；小角贝塞尔解或二维大角式都不能无条件代替这一问题。

\textbf{轴对称曲率的更精确说明}

真实硬币周围的液面是轴对称曲面，不是无限长直边旁的二维液面。若 \(x\)
表示到硬币轴线的径向距离，完整的小坡度杨---拉普拉斯方程应写成

\[
y''+\frac1x y'-\frac1{a^2}y=0.
\]

其满足远处衰减的解为修正贝塞尔函数

\[
y(x)=-l\frac{K_0(x/a)}{K_0(r/a)}.
\]

边缘斜率为

\[
y'(r)=\frac la\frac{K_1(r/a)}{K_0(r/a)}.
\]

若忽略方位曲率项
\(y'/x\)，局部二维近似给出指数解。这里必须说明近似层级：当只要求沿某一径向截面的局部恢复，或
\(r\gg a\) 使 \(1/x\)
项相对较小时，指数解合理；若要求严格轴对称解，则应使用贝塞尔函数。

\textbf{符号核验}

所得 \(y(x)<0\)、\(y'(x)>0\)、\(y''(x)<0\)。代入 \(\sigma y''=\rho_0gy\)
时，两边均为负，符号一致。若写出的指数解使
\(y''>0\)，说明曲率或静水压的符号在建方程时已经反了。

\newpage

\section{第 3 题}\label{ux7b2c-3-ux9898}

一数学摆由质量为 \(m\) 的金属小球和长度为 \(l\)
的细导线组成，在水平均匀磁场中作小幅振动。磁感应强度为
\(B\)，磁场方向垂直于摆的振动平面。摆偏离竖直线的初始最大角为
\(\alpha_0\)。当摆球通过平衡位置时，借助柔软细导线把摆线与电容为 \(C\)
的电容器短暂接通，并在很短的接触时间内使电容器完成充电。求充电后摆的最大偏角如何变化。

\begin{figure}[H]
\centering
\includegraphics[width=0.734\linewidth,height=\textheight,keepaspectratio,alt={图 3-1　磁场中的导电摆线在平衡位置短暂接入电容器。}]{../figures/question-03.pdf}
\caption*{图 3-1　磁场中的导电摆线在平衡位置短暂接入电容器。}
\end{figure}

\subsection{一、题型定位}\label{ux4e00ux9898ux578bux5b9aux4f4d-2}

导电摆线在磁场中摆动，短暂接通电容器时，充电电流受到磁力并改变摆的角动量。充电过程用角冲量关系处理，脱离电路后再用机械能求最大摆角。

\subsection{二、知识点}\label{ux4e8cux77e5ux8bc6ux70b9-2}

\begin{itemize}
\item
  转动导体上的动生电动势

  \[\mathcal E=\int(\bm v\times\bm B)\cdot \mathrm{d}\bm l.\]
\item
  变速转动时电动势与瞬时角速度的关系。
\item
  电流元在磁场中的安培力及其对悬点的力矩。
\item
  角冲量定理

  \[\Delta L_O=\int\tau_O\,\mathrm{d}t.\]
\item
  电荷量与电流的时间积分 \(q=\int I\,\mathrm{d}t\)。
\item
  摆在平衡位置和最大偏角之间的机械能守恒。
\item
  一阶近似与自洽终态方程的区别。
\end{itemize}

\subsection{三、物理图像}\label{ux4e09ux7269ux7406ux56feux50cf-2}

摆线绕悬点转动时，距悬点为 \(x\) 的线元速度为
\(\omega x\)。磁场垂直摆面，因此 \(\bm v\times\bm B\)
沿摆线方向，在悬点与摆球之间形成动生电动势。接通电容器后，电荷沿摆线转移，磁场对电流产生的力矩反抗原运动，使角速度减小。

短暂充电阶段内摆角几乎不变，可把摆线看成通过竖直位置；重力对悬点的力矩在该位置本来就是零，因此其角冲量比电磁角冲量更高阶，可以忽略。

\subsection{四、解题入口}\label{ux56dbux89e3ux9898ux5165ux53e3-2}

设接通前后平衡位置角速度分别为
\(\omega_0\)、\(\omega_1\)。依次完成三件事：

\begin{enumerate}
\def\labelenumi{\arabic{enumi}.}
\tightlist
\item
  对摆线上距悬点 \(x\) 的微元积分，得到 \(\mathcal E(\omega)\)；
\item
  把安培力矩对时间积分，用总转移电荷 \(q\) 表示角冲量；
\item
  用终态充电条件 \(q=C\mathcal E(\omega_1)\) 闭合，再由机械能关系求
  \(\alpha_1\)。
\end{enumerate}

\subsection{五、解析}\label{ux4e94ux89e3ux6790-2}

图 3-2 左侧把摆线分成距悬点为 \(x\)、长度为 \(\mathrm{d}x\) 的微元，标出
\(v(x)=\omega x\)
与微元动生电动势；右侧按``初始角速度---电容充电---终态角速度''画出能量与角冲量的传递链。第（1）步使用左图，第（2）---（4）步使用右图。

\begin{figure}[H]
\centering
\includegraphics[width=0.9\linewidth,height=\textheight,keepaspectratio,alt={图 3-2　动生电动势沿摆线积分及充电过程的角冲量关系。}]{../figures/solution-03.pdf}
\caption*{图 3-2　动生电动势沿摆线积分及充电过程的角冲量关系。}
\end{figure}

\textbf{（1）求摆线转动产生的动生电动势} 按图 3-2
左侧，摆通过平衡位置时摆线近似竖直。取距悬点为 \(x\)、长度为
\(\mathrm{d}x\) 的微元；它绕悬点转动，线速度大小为

\[v(x)=\omega x.\]

由于速度、磁场和摆线方向彼此正交，微元上的动生电动势大小为

\[\mathrm{d}\mathcal E=Bv(x)\,\mathrm{d}x=B\omega x\,\mathrm{d}x.\]

从悬点积分到摆球：

\[
\begin{aligned}
\mathcal E(\omega)
&=\int_0^lB\omega x\,\mathrm{d}x\\
&=B\omega\left[\frac{x^2}{2}\right]_0^l\\
&=\frac12Bl^2\omega.
\end{aligned}
\]

所以

\[
\boxed{\mathcal E(\omega)=\frac12Bl^2\omega}.
\]

\textbf{（2）求充电电流对摆的电磁角冲量}

图 3-2
右侧蓝、橙、绿三个框说明：电流的具体时间曲线不必先求，只需用总转移电荷
\(q\) 把磁力矩冲量连接初末角速度。设某时刻摆线中电流为 \(I(t)\)。长度为
\(\mathrm{d}x\) 的线元所受安培力大小为

\[\mathrm{d}F=IB\,\mathrm{d}x.\]

该力对悬点的力臂为 \(x\)，故微元力矩大小为

\[\mathrm{d}\tau=x\,\mathrm{d}F=IBx\,\mathrm{d}x.\]

沿摆线积分：

\[
\begin{aligned}
\tau(t)
&=\int_0^lIBx\,\mathrm{d}x\\
&=IB\left[\frac{x^2}{2}\right]_0^l\\
&=\frac12Bl^2I(t).
\end{aligned}
\]

楞次定律决定该力矩阻碍摆的运动。若以原角速度方向为正，则

\[\tau=-\frac12Bl^2I(t).\]

摆球视为质点，绕悬点的转动惯量为

\[I_O=ml^2.\]

接通时间极短，角冲量定理给出

\[
ml^2(\omega_1-\omega_0)
=\int\tau\,\mathrm{d}t.
\]

代入力矩：

\[
ml^2(\omega_1-\omega_0)
=-\frac12Bl^2\int I(t)\,\mathrm{d}t.
\]

电流的时间积分就是转移到电容器上的电荷量

\[q=\int I(t)\,\mathrm{d}t.\]

因此

\[
\boxed{
ml^2(\omega_0-\omega_1)=\frac12Bl^2q
}.
\]

\textbf{（3）使用终态充电条件自洽求 \(\omega_1\)}

这里把电容器初态取为未充电，忽略回路自感，接触时间长于充电衰减时间、短于摆的机械周期。``完成充电''按电流已经衰减理解；严格零电阻且保留自感的回路会持续振荡，不能直接使用本题的终态电压条件。

充电结束时电流为零，电容器电压与此时摆线的动生电动势相等。由于角速度已在充电过程中由
\(\omega_0\) 降到 \(\omega_1\)，终态应写成

\[
\frac qC=\mathcal E(\omega_1)
=\frac12Bl^2\omega_1.
\]

于是

\[
\boxed{q=\frac12CBl^2\omega_1}.
\]

把它代入角冲量式：

\[
ml^2(\omega_0-\omega_1)
=\frac12Bl^2\left(\frac12CBl^2\omega_1\right).
\]

右端逐项相乘：

\[
ml^2(\omega_0-\omega_1)
=\frac14CB^2l^4\omega_1.
\]

两边除以 \(ml^2\)：

\[
\omega_0-\omega_1
=\frac{CB^2l^2}{4m}\omega_1.
\]

定义无量纲耦合参数

\[
\beta=\frac{CB^2l^2}{4m}.
\]

则

\[\omega_0-\omega_1=\beta\omega_1,\]

从而

\[\omega_0=(1+\beta)\omega_1.\]

所以

\[
\boxed{
\omega_1=\frac{\omega_0}{1+\beta}
=\frac{\omega_0}{1+CB^2l^2/(4m)}
}.
\]

\textbf{（4）由平衡位置角速度求最大偏角}

接通前，摆从最大偏角 \(\alpha_0\) 摆到最低点，机械能守恒：

\[
mgl(1-\cos\alpha_0)=\frac12ml^2\omega_0^2.
\]

接触结束后，摆以角速度 \(\omega_1\) 从最低点上升到新的最大偏角
\(\alpha_1\)：

\[
\frac12ml^2\omega_1^2=mgl(1-\cos\alpha_1).
\]

第二式除以第一式，\(m\)、\(g\)、\(l\) 全部约去：

\[
\frac{1-\cos\alpha_1}{1-\cos\alpha_0}
=\frac{\omega_1^2}{\omega_0^2}.
\]

代入 \(\omega_1/\omega_0=1/(1+\beta)\)：

\[
\boxed{
1-\cos\alpha_1
=\frac{1-\cos\alpha_0}
{\left(1+\dfrac{CB^2l^2}{4m}\right)^2}
}.
\]

若摆角确实很小，使用

\[1-\cos\alpha\simeq\frac12\alpha^2,\]

得到更直观的振幅关系

\[
\boxed{
\alpha_1\simeq
\frac{\alpha_0}{1+CB^2l^2/(4m)}
}.
\]

\textbf{一阶近似与自洽终态方程的关系}

在一阶近似下，可把电荷量写成

\[q\simeq C\mathcal E(\omega_0)
=\frac12CBl^2\omega_0,\]

因而得到

\[\omega_1\simeq(1-\beta)\omega_0\]

以及

\[1-\cos\alpha_1\simeq(1-\beta)^2(1-\cos\alpha_0).\]

这个写法把充电过程中角速度的下降忽略在电动势计算之外，只能看作
\(\beta\ll1\) 的一阶近似。因为

\[
\frac1{1+\beta}=1-\beta+\beta^2-\cdots,
\]

两种结果在一阶确实一致。若 \(\beta\) 不小，\((1-\beta)\omega_0\)
甚至会在 \(\beta>1\)
时给出反向角速度，显然不能作为一般终态。因此教师版采用终态电压与终态动生电动势相等的自洽结果。

\textbf{能量去向与串联电阻核验}

接通过程并非摆的机械能守恒过程。机械能的一部分储存在电容器中，实际柔性导线与瞬态电流还可能产生焦耳热和电磁辐射。若以串联电阻
\(R\) 完成充电，记 \(a_B=Bl^2/2\)、\(J=ml^2\)，由
\(J\omega+a_Bq=J\omega_0\) 及 \(R\dot q+q/C=a_B\omega\)，得到

\[R\dot q+\frac{1+\beta}{C}q=a_B\omega_0.\]

衰减时间为 \(RC/(1+\beta)\)。设充电前动能为 \(K_0=J\omega_0^2/2\)，则

\[K_1=\frac{K_0}{(1+\beta)^2},\quad
U_C=\frac{K_0\beta}{(1+\beta)^2},\quad
Q_R=\frac{K_0\beta}{1+\beta}.\]

三者之和为
\(K_0\)。因此充电后的幅度计算与总能量收支相容，但机械能损失大于电容储能。

\newpage

\section{第 4 题}\label{ux7b2c-4-ux9898}

一种发光体服从余弦发光定律。对面积为 \(\mathrm{d}S\) 的发光面元，在时间
\(\mathrm{d}t\) 内向立体角 \(\mathrm{d}\Omega\) 中发出的能量为

\[
\mathrm{d}E=I_0\cos\theta\,\mathrm{d}\Omega\,\mathrm{d}S\,\mathrm{d}t,
\]

其中 \(\theta\) 为发光方向与面元法线的夹角，\(I_0\) 为常量。

\begin{enumerate}
\def\labelenumi{\arabic{enumi}.}
\tightlist
\item
  求单位发光面积在单位时间内发出的总功率 \(P_0\)。
\item
  一个无限大平面由上述面元构成。在距它 \(l_0\)
  处放置一个与之平行的小接收面，求接收面单位面积得到的功率 \(P_1\)。
\item
  上述发光面构成一根足够长、半径为 \(r\) 的圆柱。在距圆柱轴线 \(x\)
  处放置小接收面，并可任意选择其朝向。求单位接收面积所得最大功率
  \(P_2(x)\)，取 \(x>r\)。
\item
  两根足够长、半径均为 \(r\) 的圆柱灯安装在天花板上，轴线均沿 \(y\)
  方向，横坐标分别为 \(x=\pm h\)，且 \(h\gg r\)。桌面与天花板的距离为
  \(\lambda h\)。要求桌面中心 \(x=0\) 附近的单位面积受光功率随 \(x\)
  的变化尽可能平坦，求 \(\lambda\)。
\end{enumerate}

\begin{figure}[H]
\centering
\includegraphics[width=0.811\linewidth,height=\textheight,keepaspectratio,alt={图 4-1　余弦发光面元的法线、发光方向与微小立体角。}]{../figures/question-04.pdf}
\caption*{图 4-1　余弦发光面元的法线、发光方向与微小立体角。}
\end{figure}

\subsection{一、题型定位}\label{ux4e00ux9898ux578bux5b9aux4f4d-3}

先积分单个朗伯面元对接收面的照度贡献，再叠加对称灯源。中心照度要尽量平坦，可把二阶空间变化作为求灯距的条件。

\subsection{二、知识点}\label{ux4e8cux77e5ux8bc6ux70b9-3}

\begin{itemize}
\item
  立体角元
  \(\mathrm{d}\Omega=\sin\theta\,\mathrm{d}\theta\,\mathrm{d}\varphi\)。
\item
  朗伯面元总辐射功率的半球积分。
\item
  发射端和接收端各自的余弦投影因子。
\item
  面元间传输关系

  \[\mathrm{d}P=I_0\frac{\cos\theta_s\cos\theta_r}{R^2}\,\mathrm{d}S_s\,\mathrm{d}S_r.\]
\item
  无限平面的极坐标积分。
\item
  长圆柱对外部点的投影立体角与 \(1/x\) 规律。
\item
  对称函数在原点的一阶导数自动为零，平坦化条件为二阶导数为零。
\end{itemize}

\subsection{三、物理图像}\label{ux4e09ux7269ux7406ux56feux50cf-3}

余弦定律中的 \(\cos\theta\)
表示同一面元向斜方向的辐射变弱。接收面还会因斜入射产生另一个投影因子。无限发光平面虽然远处面元更多，但距离衰减与投影因子恰好使总照度与高度
\(l_0\) 无关。

两根灯关于 \(x=0\)
对称，所以中心处的线性变化自动抵消。要让中心附近``尽可能平坦''，就要进一步消去二次项；此时四次项成为首个非恒定项。

\subsection{四、解题入口}\label{ux56dbux89e3ux9898ux5165ux53e3-3}

第一问直接在发光面元上方半球积分。第二问以接收面正下方为极点，令发光面上的环带半径为
\(s\)。第三问可把无限长圆柱视为在接收面视野中的一条带状辐射源，利用投影立体角积分。第四问写出两灯贡献之和并令
\(P''(0)=0\)。

\subsection{五、解析}\label{ux4e94ux89e3ux6790-3}

图 4-2 左侧把无限发光平面分成半径 \(s\)、宽度 \(\mathrm{d}s\)
的环带，并标出接收面高度 \(l_0\)；右侧画出两根圆柱灯到桌面点 \(x\)
的距离和投影几何。第（2）步使用左图，第（4）---（5）步使用右图；第（1）、（3）步先建立后续所需的辐射定律。

\begin{figure}[H]
\centering
\includegraphics[width=0.9\linewidth,height=\textheight,keepaspectratio,alt={图 4-2　无限平面照度积分及双圆柱灯的中心平坦化几何。}]{../figures/solution-04.pdf}
\caption*{图 4-2　无限平面照度积分及双圆柱灯的中心平坦化几何。}
\end{figure}

\textbf{（1）单位发光面积的总功率} 由定义，单位面积、单位时间向
\(\mathrm{d}\Omega\) 发出的功率为

\[\mathrm{d}P=I_0\cos\theta\,\mathrm{d}\Omega.\]

发光只在面元外侧半球内进行，所以

\[0\le\theta\le\frac\pi2,
\qquad
0\le\varphi<2\pi.\]

球坐标中的立体角元为

\[\mathrm{d}\Omega=\sin\theta\,\mathrm{d}\theta\,\mathrm{d}\varphi.\]

于是

\[
\begin{aligned}
P_0
&=\int_0^{2\pi}\int_0^{\pi/2}
I_0\cos\theta\sin\theta\,\mathrm{d}\theta\,\mathrm{d}\varphi\\
&=I_0\int_0^{2\pi}\mathrm{d}\varphi
\int_0^{\pi/2}\cos\theta\sin\theta\,\mathrm{d}\theta\\
&=I_0(2\pi)
\left[\frac12\sin^2\theta\right]_0^{\pi/2}\\
&=I_0(2\pi)\cdot\frac12\\
&=\pi I_0.
\end{aligned}
\]

因此

\[\boxed{P_0=\pi I_0}.\]

\textbf{（2）无限发光平面对平行小面的照度}

按图 4-2 左侧，在接收面上取单位面积
\(\mathrm{d}S_r=1\)。以其在发光平面上的垂足为中心，取半径为
\(s\)、宽度为 \(\mathrm{d}s\) 的环带。环带面积为

\[\mathrm{d}S_s=2\pi s\,\mathrm{d}s.\]

环带到接收点的距离为

\[R=\sqrt{s^2+l_0^2}.\]

发光面法线和接收面法线都与两点连线具有相同夹角，故

\[
\cos\theta_s=\cos\theta_r=\frac{l_0}{R}.
\]

发光环带在接收单位面积处张成的立体角为

\[
\mathrm{d}\Omega=\frac{\cos\theta_r\,\mathrm{d}S_r}{R^2}
=\frac{l_0}{R^3}.
\]

环带对接收面的功率贡献为

\[
\begin{aligned}
\mathrm{d}P_1
&=I_0\cos\theta_s\,\mathrm{d}S_s\,\mathrm{d}\Omega\\
&=I_0\frac{l_0}{R}(2\pi s\,\mathrm{d}s)
\frac{l_0}{R^3}\\
&=\frac{2\pi I_0l_0^2s\,\mathrm{d}s}
{(s^2+l_0^2)^2}.
\end{aligned}
\]

从 \(s=0\) 积分到无穷远：

\[
P_1=2\pi I_0l_0^2
\int_0^\infty\frac{s\,\mathrm{d}s}{(s^2+l_0^2)^2}.
\]

令

\[u=s^2+l_0^2,
\qquad \mathrm{d}u=2s\,\mathrm{d}s,\]

则

\[
\begin{aligned}
\int_0^\infty\frac{s\,\mathrm{d}s}{(s^2+l_0^2)^2}
&=\frac12\int_{l_0^2}^{\infty}u^{-2}\,\mathrm{d}u\\
&=\frac12\left[-u^{-1}\right]_{l_0^2}^{\infty}\\
&=\frac1{2l_0^2}.
\end{aligned}
\]

因此

\[
P_1=2\pi I_0l_0^2\cdot\frac1{2l_0^2}
=\pi I_0.
\]

即

\[\boxed{P_1=\pi I_0=P_0}.\]

结果与 \(l_0\)
无关。高度增大使单个面元贡献减小，但可见发光面积按同一比例扩大。

\textbf{（3）无限长圆柱灯的最大照度}

接收面朝向可调。由于圆柱轴向无限长且几何关于``接收点---圆柱轴''所在平面对称，最大照度时接收面法线应指向圆柱轴线。若法线偏离该方向，所有入射光在接收面上的投影总量都会减小。

下面把投影立体角结论完整积分出来。令圆柱轴为 \(z\) 轴，接收点为
\(P=(x,0,0)\)，其中 \(x>r\)。圆柱面元取为

\[
S=(r\cos\varphi,r\sin\varphi,z),
\qquad
\mathrm{d}S=r\,\mathrm{d}\varphi\,\mathrm{d}z.
\]

从面元指向接收点的矢量为

\[
\bm R=(x-r\cos\varphi,-r\sin\varphi,-z),
\]

其长度满足

\[
R^2=x^2+r^2-2xr\cos\varphi+z^2.
\]

圆柱面外法线为

\[\hat{\bm n}_s=(\cos\varphi,\sin\varphi,0).\]

故发射余弦为

\[
\cos\theta_s
=\hat{\bm n}_s\cdot\frac{\bm R}{R}
=\frac{x\cos\varphi-r}{R}.
\]

只有 \(x\cos\varphi-r>0\) 的圆柱面可见。令

\[\varphi_0=\arccos\frac rx,\]

则可见范围为

\[-\varphi_0<\varphi<\varphi_0.\]

接收面法线指向圆柱轴，即

\[\hat{\bm n}_r=(-1,0,0).\]

从接收点指向面元的方向为 \(-\bm R/R\)，所以接收余弦为

\[
\cos\theta_r
=\hat{\bm n}_r\cdot\frac{-\bm R}{R}
=\frac{x-r\cos\varphi}{R}.
\]

单位接收面积所得微元功率为

\[
\mathrm{d}P=I_0\cos\theta_s\cos\theta_r
\frac{\mathrm{d}S}{R^2}.
\]

代入各量：

\[
\mathrm{d}P=I_0r
\frac{(x\cos\varphi-r)(x-r\cos\varphi)}
{[x^2+r^2-2xr\cos\varphi+z^2]^2}
\,\mathrm{d}\varphi\,\mathrm{d}z.
\]

记

\[a^2=x^2+r^2-2xr\cos\varphi.\]

先沿无限长轴向积分，并使用

\[
\int_{-\infty}^{\infty}
\frac{\mathrm{d}z}{(z^2+a^2)^2}
=\frac{\pi}{2a^3}.
\]

于是

\[
P_2=\frac{\pi I_0r}{2}
\int_{-\varphi_0}^{\varphi_0}
\frac{(x\cos\varphi-r)(x-r\cos\varphi)}
{(x^2+r^2-2xr\cos\varphi)^{3/2}}\,\mathrm{d}\varphi.
\]

被积函数恰好是下式的导数：

\[
\frac{\mathrm{d}}{\mathrm{d}\varphi}
\left[
\frac{\sin\varphi}
{\sqrt{x^2+r^2-2xr\cos\varphi}}
\right]
=\frac{(x\cos\varphi-r)(x-r\cos\varphi)}
{(x^2+r^2-2xr\cos\varphi)^{3/2}}.
\]

因此角积分等于

\[
\left.
\frac{\sin\varphi}
{\sqrt{x^2+r^2-2xr\cos\varphi}}
\right|_{-\varphi_0}^{\varphi_0}.
\]

在边界 \(\cos\varphi_0=r/x\) 处，

\[
\sin\varphi_0=\frac{\sqrt{x^2-r^2}}{x},
\]

且

\[
\sqrt{x^2+r^2-2xr\cos\varphi_0}
=\sqrt{x^2-r^2}.
\]

故上限的比值为 \(1/x\)，下限为 \(-1/x\)，整个角积分为 \(2/x\)。最终

\[
P_2=\frac{\pi I_0r}{2}\cdot\frac2x
=\pi I_0\frac rx.
\]

所以

\[
\boxed{P_2(x)=\pi I_0\frac rx},
\qquad x>r.
\]

两个极限可作检查：当 \(x\gg r\) 时，照度按 \(1/x\)
衰减，这是无限长线光源的规律；当 \(x\to r^+\)
时，\(P_2\to\pi I_0=P_0\)，接收面贴近圆柱且正对局部表面时回到无限朗伯面的结果。

\textbf{（4）两根圆柱灯在桌面上的照度函数}

转看图 4-2 右侧。两灯轴间距为 \(2h\)，桌面到灯轴高度为
\(\lambda h\)；取桌面上一点的横坐标为
\(x\)，它到左、右两根灯轴的距离分别为

\[
d_- =\sqrt{(x+h)^2+(\lambda h)^2},
\]

\[
d_+ =\sqrt{(x-h)^2+(\lambda h)^2}.
\]

桌面法线竖直向上，而每根灯的最佳接收方向是沿该点指向灯轴。由于桌面不能转动，还要乘以入射方向与桌面法线的余弦：

\[
\cos\gamma_\pm=\frac{\lambda h}{d_\pm}.
\]

第 3 问给出的最佳定向照度为
\(\pi I_0r/d_\pm\)，投影到水平桌面后，每根灯贡献

\[
P_\pm(x)=\frac{\pi I_0r}{d_\pm}
\frac{\lambda h}{d_\pm}
=\frac{\pi I_0r\lambda h}{d_\pm^2}.
\]

总照度为

\[
\boxed{
P(x)=\pi I_0r\lambda h
\left[
\frac1{(x+h)^2+\lambda^2h^2}
+\frac1{(x-h)^2+\lambda^2h^2}
\right]
}.
\]

\textbf{（5）在 \(x=0\) 附近作二阶展开}

令

\[u=\frac xh,
\qquad A=1+\lambda^2.\]

则除去与 \(u\) 无关的正系数后，需要研究

\[
f(u)=\frac1{A+2u+u^2}
+\frac1{A-2u+u^2}.
\]

对第一项使用

\[
\frac1{A+\varepsilon}
=\frac1A-\frac{\varepsilon}{A^2}
+\frac{\varepsilon^2}{A^3}+O(\varepsilon^3).
\]

保留到 \(u^2\)：

\[
\frac1{A\pm2u+u^2}
=\frac1A
\mp\frac{2u}{A^2}
-\frac{u^2}{A^2}
+\frac{4u^2}{A^3}
+O(u^3).
\]

两式相加，奇次项抵消：

\[
\begin{aligned}
f(u)
&=\frac2A
+2\left(-\frac1{A^2}+\frac4{A^3}\right)u^2
+O(u^4)\\
&=\frac2A
+\frac{2(4-A)}{A^3}u^2
+O(u^4)\\
&=\frac2{1+\lambda^2}
+\frac{2(3-\lambda^2)}{(1+\lambda^2)^3}u^2
+O(u^4).
\end{aligned}
\]

对称性已经保证一次项为零。要使中心附近进一步平坦，令二次项系数为零：

\[3-\lambda^2=0.\]

因 \(\lambda>0\)，故

\[\boxed{\lambda=\sqrt3}.\]

此时中心附近的首个非恒定修正是 \(x^4\) 量级。

\textbf{投影立体角的等价理解}

上面的双重积分也可概括为``照度等于辐亮度 \(I_0\)
乘以光源的投影立体角''。圆柱沿轴向无限延伸，使投影立体角按 \(r/x\) 而非
\(r^2/x^2\) 缩放。教师可先用完整积分保证因子无误，再用投影立体角解释
\(1/x\) 的物理来源。

\newpage

\section{第 5 题}\label{ux7b2c-5-ux9898}

两个绝热刚性容器的体积均为
\(V_0\)，中间由一根绝热细管相连，管内装有一台理想、缓慢转动的绝热汽轮机。初始时，左容器装有单原子理想气体，压强为
\(p_0\)、温度为 \(T_0\)，定容摩尔热容为

\[C_V=\frac32R,\]

右容器为真空。汽轮机极缓慢地对外做功；气体从左侧流出时可视为经历准静态绝热膨胀，随后进入右容器并均匀混合。气体膨胀所做的功全部转化为汽轮机对外输出的功。

\begin{enumerate}
\def\labelenumi{\arabic{enumi}.}
\tightlist
\item
  当左容器中气体压强降为 \(\lambda p_0\) 时，求左侧温度
  \(T_L(\lambda)\)。
\item
  在同一状态下，只对全部气体使用热力学第一定律。若汽轮机累计做功为
  \(W\)，求右侧温度 \(T_R(\lambda)\) 和压强 \(p_R(\lambda)\)。
\item
  对通过汽轮机的一小团气体使用热力学第一定律，建立 \(W\) 关于
  \(\lambda\) 的微分方程，并求出状态函数。
\end{enumerate}

进一步求系统达到力学平衡时汽轮机的总功 \(W_m\) 及气体的总熵变。

\begin{figure}[H]
\centering
\includegraphics[width=0.9\linewidth,height=\textheight,keepaspectratio,alt={图 5-1　两个绝热刚性容器、缓慢绝热汽轮机与气体迁移过程。}]{../figures/question-05.pdf}
\caption*{图 5-1　两个绝热刚性容器、缓慢绝热汽轮机与气体迁移过程。}
\end{figure}

\subsection{一、题型定位}\label{ux4e00ux9898ux578bux5b9aux4f4d-4}

两侧绝热容器通过汽轮机交换气体和能量。容器虽刚性，流出气体仍携带焓；对整个系统列能量账本才能确定最终状态。

\subsection{二、知识点}\label{ux4e8cux77e5ux8bc6ux70b9-4}

\begin{itemize}
\item
  开口系统流出物质携带的能量是摩尔焓 \(h=C_PT\)。
\item
  单原子理想气体

  \[C_V=\frac32R,
  \qquad C_P=\frac52R,
  \qquad \gamma=\frac53.\]
\item
  刚性绝热容器的放气关系
  \(T\propto n^{\gamma-1}\)、\(p\propto n^\gamma\)。
\item
  整体第一定律 \(U_0=U_L+U_R+W\)。
\item
  理想绝热汽轮机的等熵关系

  \[\frac{T_e}{T_L}=\left(\frac{p_R}{p_L}\right)^{(\gamma-1)/\gamma}.\]
\item
  无量纲化与可分离变量微分方程。
\item
  力学平衡条件 \(p_L=p_R\)。
\item
  理想气体熵变及可逆过程的熵守恒。
\end{itemize}

\subsection{三、物理图像}\label{ux4e09ux7269ux7406ux56feux50cf-4}

左容器的气体不断减少。每流出一小团气体，左侧剩余气体要对它做流动功，所以左侧降温速度比简单``按内部能带走''更快。气体经过理想透平继续绝热膨胀，焓降低量转化为轴功，然后进入右容器。

由于过程足够缓慢，透平出口压强可视为等于当时右容器压强。理想等熵膨胀还会使出口温度恰好与右侧均匀气体温度相容，因此理想极限下不会因温差混合额外产生熵。

\subsection{四、解题入口}\label{ux56dbux89e3ux9898ux5165ux53e3-4}

令初始物质的量为

\[n_0=\frac{p_0V_0}{RT_0},\]

左、右物质的量分别为
\(n_L\)、\(n_R=n_0-n_L\)。第一问先从左容器开口能量方程求出
\(T_L(n_L)\)，再结合 \(p_LV_0=n_LRT_L\) 把 \(n_L\) 换成 \(\lambda\)。

第二问引入

\[w=\frac{W}{U_0}=\frac{2W}{3p_0V_0},\]

可显著缩短公式。第三问再令

\[y=\frac{p_R}{p_0}=1-\lambda-w,\]

微分方程会化成可直接分离变量的形式。

\subsection{五、解析}\label{ux4e94ux89e3ux6790-4}

图 5-2 左侧画出刚性左室放气时 \(T_L/T_0=\lambda^{2/5}\)
的状态曲线，右侧把一微元气体从左室经绝热透平进入右室的能量流程画成方框图。第（1）步对应左图，第（2）---（4）步沿右图从整体守恒推进到微分方程。

\begin{figure}[H]
\centering
\includegraphics[width=0.9\linewidth,height=\textheight,keepaspectratio,alt={图 5-2　左室绝热放气曲线与微元气体经透平进入右室的能量流程。}]{../figures/solution-05.pdf}
\caption*{图 5-2　左室绝热放气曲线与微元气体经透平进入右室的能量流程。}
\end{figure}

\textbf{（1）左容器的绝热放气关系} 图 5-2
左图的曲线不是封闭气体的普通绝热线，而是开口刚性容器放气曲线。左容器内能为

\[U_L=n_LC_VT_L.\]

设左侧物质的量发生微小变化 \(\mathrm{d}n_L<0\)。流出的物质的量为
\(-\mathrm{d}n_L>0\)。开口系统的绝热能量平衡为：左容器内能的减少等于流出气体携带走的焓。因此

\[
\mathrm{d}U_L=C_PT_L\,\mathrm{d}n_L.
\]

这里右端为负，因为 \(\mathrm{d}n_L<0\)。展开左端：

\[
d(n_LC_VT_L)
=C_Vn_L\,\mathrm{d}T_L+C_VT_L\,\mathrm{d}n_L.
\]

所以

\[
C_Vn_L\,\mathrm{d}T_L+C_VT_L\,\mathrm{d}n_L
=C_PT_L\,\mathrm{d}n_L.
\]

把含 \(\mathrm{d}n_L\) 的项合并：

\[
C_Vn_L\,\mathrm{d}T_L
=(C_P-C_V)T_L\,\mathrm{d}n_L.
\]

利用迈耶关系 \(C_P-C_V=R\)：

\[
C_Vn_L\,\mathrm{d}T_L=RT_L\,\mathrm{d}n_L.
\]

两边除以 \(C_Vn_LT_L\)：

\[
\frac{\mathrm{d}T_L}{T_L}
=\frac{R}{C_V}\frac{\mathrm{d}n_L}{n_L}.
\]

单原子理想气体有 \(R/C_V=2/3\)，故

\[
\frac{\mathrm{d}T_L}{T_L}=\frac23\frac{\mathrm{d}n_L}{n_L}.
\]

从初态 \((n_0,T_0)\) 积分到 \((n_L,T_L)\)：

\[
\int_{T_0}^{T_L}\frac{\mathrm{d}T}{T}
=\frac23\int_{n_0}^{n_L}\frac{\mathrm{d}n}{n}.
\]

于是

\[
\ln\frac{T_L}{T_0}
=\frac23\ln\frac{n_L}{n_0},
\]

从而

\[
\frac{T_L}{T_0}
=\left(\frac{n_L}{n_0}\right)^{2/3}.
\]

另一方面，状态方程给出

\[
\frac{p_L}{p_0}
=\frac{n_LT_L}{n_0T_0}
=\left(\frac{n_L}{n_0}\right)^{1+2/3}
=\left(\frac{n_L}{n_0}\right)^{5/3}.
\]

题设 \(p_L=\lambda p_0\)，所以

\[
\lambda=\left(\frac{n_L}{n_0}\right)^{5/3}.
\]

两边取 \(3/5\) 次幂：

\[
\boxed{\frac{n_L}{n_0}=\lambda^{3/5}}.
\]

代回温度关系：

\[
\frac{T_L}{T_0}
=\left(\lambda^{3/5}\right)^{2/3}
=\lambda^{2/5}.
\]

故

\[\boxed{T_L(\lambda)=T_0\lambda^{2/5}}.\]

\textbf{（2）只用整体第一定律求右侧状态与 \(W\) 的关系}

按图 5-2 右侧的整体能量流程，透平输出功 \(W\)
从两室总内能中扣除。初态全部气体内能为

\[
U_0=n_0C_VT_0
=\frac32n_0RT_0
=\frac32p_0V_0.
\]

左容器当前内能可直接由 \(p_LV_0=n_LRT_L\) 写成

\[
U_L=\frac32n_LRT_L
=\frac32p_LV_0
=\frac32\lambda p_0V_0.
\]

右容器内能为

\[U_R=\frac32p_RV_0.\]

两个容器和管道均绝热，系统对外只输出汽轮机功 \(W\)，因此整体第一定律为

\[U_0=U_L+U_R+W.\]

代入：

\[
\frac32p_0V_0
=\frac32\lambda p_0V_0
+\frac32p_RV_0+W.
\]

两边乘 \(2/(3V_0)\)：

\[p_0=\lambda p_0+p_R+\frac{2W}{3V_0}.\]

所以

\[
\boxed{
p_R=(1-\lambda)p_0-\frac{2W}{3V_0}
}.
\]

定义

\[w=\frac{2W}{3p_0V_0},\]

则

\[
\boxed{\frac{p_R}{p_0}=1-\lambda-w}.
\]

右侧物质的量为

\[
n_R=n_0-n_L
=n_0(1-\lambda^{3/5}).
\]

利用右侧状态方程

\[p_RV_0=n_RRT_R\]

与 \(p_0V_0=n_0RT_0\)，两式相除：

\[
\frac{p_R}{p_0}
=\frac{n_R}{n_0}\frac{T_R}{T_0}.
\]

因此

\[
\boxed{
T_R=T_0\frac{1-\lambda-w}{1-\lambda^{3/5}}
}.
\]

写回 \(W\)：

\[
\boxed{
T_R=T_0
\frac{1-\lambda-\dfrac{2W}{3p_0V_0}}
{1-\lambda^{3/5}}
}.
\]

到这里仅用了整体第一定律，所以 \(W\) 仍然未知。

\textbf{（3）对通过透平的微元气体列式}

现在沿图 5-2 右侧箭头跟踪一小团气体：它从左侧状态 \((p_L,T_L)\)
出发，经透平等熵膨胀，再进入右侧状态 \((p_R,T_R)\)。左侧压强为

\[p_L=\lambda p_0,\]

右侧压强记为

\[p_R=yp_0,\]

即

\[y=1-\lambda-w.\]

一小团气体从左侧状态 \((p_L,T_L)\) 经理想缓慢绝热透平膨胀到出口压强
\(p_R\)。理想绝热且准静态意味着等熵，故

\[
\frac{T_e}{T_L}
=\left(\frac{p_R}{p_L}\right)^{(\gamma-1)/\gamma}.
\]

对 \(\gamma=5/3\)，指数为

\[
\frac{\gamma-1}{\gamma}
=\frac{2/3}{5/3}
=\frac25.
\]

于是

\[
T_e=T_L\left(\frac y\lambda\right)^{2/5}.
\]

代入 \(T_L=T_0\lambda^{2/5}\)：

\[
\begin{aligned}
T_e
&=T_0\lambda^{2/5}
\left(\frac y\lambda\right)^{2/5}\\
&=T_0y^{2/5}.
\end{aligned}
\]

因此

\[\boxed{T_e=T_0y^{2/5}}.\]

透平单位摩尔输出功等于焓降：

\[\mathrm{d}w_{\rm mol}=C_P(T_L-T_e).\]

左侧物质的量

\[n_L=n_0\lambda^{3/5},\]

故

\[
\mathrm{d}n_L=\frac35n_0\lambda^{-2/5}\,\mathrm{d}\lambda.
\]

当 \(\mathrm{d}\lambda<0\) 时，流过透平的正物质的量为

\[
\mathrm{d}n=-\mathrm{d}n_L
=-\frac35n_0\lambda^{-2/5}\,\mathrm{d}\lambda.
\]

于是累计功的微分为

\[
\mathrm{d}W=C_P(T_L-T_e)\,\mathrm{d}n.
\]

代入 \(C_P=5R/2\)、\(T_L=T_0\lambda^{2/5}\)、\(T_e=T_0y^{2/5}\)：

\[
\begin{aligned}
\mathrm{d}W
&=\frac52R T_0
\left(\lambda^{2/5}-y^{2/5}\right)
\left(-\frac35n_0\lambda^{-2/5}\,\mathrm{d}\lambda\right)\\
&=-\frac32n_0RT_0
\left[1-\left(\frac y\lambda\right)^{2/5}\right]\mathrm{d}\lambda\\
&=-\frac32p_0V_0
\left[1-\left(\frac y\lambda\right)^{2/5}\right]\mathrm{d}\lambda.
\end{aligned}
\]

因此 \(W\) 的微分方程为

\[
\boxed{
\frac{\mathrm{d}W}{\mathrm{d}\lambda}
=-\frac32p_0V_0
\left[
1-\left(
\frac{1-\lambda-\dfrac{2W}{3p_0V_0}}{\lambda}
\right)^{2/5}
\right]
}.
\]

用 \(w=2W/(3p_0V_0)\) 表示：

\[
\boxed{
\frac{\mathrm{d}w}{\mathrm{d}\lambda}
=-1+\left(\frac{1-\lambda-w}{\lambda}\right)^{2/5}
}.
\]

初始条件为

\[w(1)=0.\]

初态右室没有气体，\(T_R\) 本身未定义；后文温度只用于
\(\lambda<1\)。理想输运解的过程范围是
\(2^{-5/3}\le\lambda<1\)，超过压强平衡后不能再沿原方向输出正功。

\textbf{（4）把微分方程化为右侧压强方程并积分}

定义

\[y=1-\lambda-w.\]

对 \(\lambda\) 求导：

\[
\frac{\mathrm{d}y}{\mathrm{d}\lambda}
=-1-\frac{\mathrm{d}w}{\mathrm{d}\lambda}.
\]

代入上式：

\[
\begin{aligned}
\frac{\mathrm{d}y}{\mathrm{d}\lambda}
&=-1-\left[-1+\left(\frac y\lambda\right)^{2/5}\right]\\
&=-\left(\frac y\lambda\right)^{2/5}.
\end{aligned}
\]

所以

\[
y^{-2/5}\,\mathrm{d}y=-\lambda^{-2/5}\,\mathrm{d}\lambda.
\]

从初态 \(\lambda=1\)、\(y=0\) 积分到当前状态：

\[
\int_0^y u^{-2/5}\,\mathrm{d}u
=-\int_1^\lambda s^{-2/5}\,\mathrm{d}s.
\]

两边积分：

\[
\frac53y^{3/5}
=-\frac53(\lambda^{3/5}-1).
\]

约去 \(5/3\)：

\[
y^{3/5}=1-\lambda^{3/5}.
\]

因此

\[
\boxed{
y=\frac{p_R}{p_0}
=\left(1-\lambda^{3/5}\right)^{5/3}
}.
\]

由 \(w=1-\lambda-y\)：

\[
\boxed{
w=1-\lambda-\left(1-\lambda^{3/5}\right)^{5/3}
}.
\]

所以

\[
\boxed{
W(\lambda)=\frac32p_0V_0
\left[
1-\lambda-\left(1-\lambda^{3/5}\right)^{5/3}
\right]
}.
\]

右侧温度也可化简。由

\[
\frac{T_R}{T_0}
=\frac{y}{1-\lambda^{3/5}},
\]

代入 \(y=(1-\lambda^{3/5})^{5/3}\)：

\[
\boxed{
T_R=T_0\left(1-\lambda^{3/5}\right)^{2/3}
}.
\]

它恰好等于透平出口温度
\(T_e=T_0y^{2/5}\)，说明理想极限下进入右室的微元与右室现有气体同温，不产生有限温差混合熵。

\textbf{（5）达到力学平衡时的总功}

力学平衡要求两侧压强相等：

\[p_L=p_R.\]

即

\[
\lambda=\left(1-\lambda^{3/5}\right)^{5/3}.
\]

两边取 \(3/5\) 次幂：

\[
\lambda^{3/5}=1-\lambda^{3/5}.
\]

所以

\[
2\lambda^{3/5}=1,
\]

\[
\lambda^{3/5}=\frac12,
\]

从而

\[
\boxed{\lambda_m=2^{-5/3}}.
\]

此时

\[
\frac{T_L}{T_0}=\lambda_m^{2/5}
=2^{-2/3},
\]

右侧温度相同：

\[
\frac{T_R}{T_0}
=\left(1-\frac12\right)^{2/3}
=2^{-2/3}.
\]

平衡时 \(y=\lambda_m\)，故

\[
w_m=1-\lambda_m-y_m
=1-2\lambda_m
=1-2^{-2/3}.
\]

于是

\[
\boxed{
W_m=\frac32p_0V_0
\left(1-2^{-2/3}\right)
}.
\]

数值上

\[
2^{-2/3}\simeq0.62996,
\]

所以

\[W_m\simeq0.55506\,p_0V_0.\]

\textbf{（6）总熵变}

平衡时左右各有一半物质：

\[n_L=n_R=\frac{n_0}{2}.\]

每一半气体的体积均为
\(V_0\)，所以相对于初态单位物质所占体积扩大为原来的两倍；温度降为
\(2^{-2/3}T_0\)。理想气体单位摩尔熵变为

\[
\Delta s=C_V\ln\frac{T_f}{T_0}
+R\ln\frac{v_f}{v_0}.
\]

代入

\[
\frac{T_f}{T_0}=2^{-2/3},
\qquad
\frac{v_f}{v_0}=2:
\]

\[
\begin{aligned}
\Delta s
&=\frac32R\ln(2^{-2/3})+R\ln2\\
&=\frac32R\left(-\frac23\ln2\right)+R\ln2\\
&=-R\ln2+R\ln2\\
&=0.
\end{aligned}
\]

因此全部气体的总熵变为

\[\boxed{\Delta S=0}.\]

\textbf{微分式中的流动功}

若只按内能差尝试写成

\[\mathrm{d}W\propto-C_V(T_L-T_R)\,\mathrm{d}n_L.\]

这一式忽略了流动气体携带的 \(pV\) 流动功。透平的轴功由焓降
\(C_P(T_L-T_e)\) 决定，而不是由内能差 \(C_V(T_L-T_R)\)
决定。若把这一内能差形式与整体能量关系联立，在 \(\lambda\to1\)
的真空初始极限会出现不自洽的首阶系数，正是漏掉流动功的信号。

此外，\(n_L=n_0\lambda^{3/5}\) 求导必须得到

\[\mathrm{d}n_L=\frac35n_0\lambda^{-2/5}\,\mathrm{d}\lambda.\]

微分方程中还必须保留这个
\(\lambda^{-2/5}\)。研究一小团通过透平的气体仍是正确入口，但必须使用焓降和等熵关系完成计算。

\textbf{为什么最后熵变为零}

题设的理想、缓慢、绝热意味着透平内等熵；解出的 \(T_e=T_R\)
又消除了有限温差混合。因而整个极限过程可逆。若透平有不可逆损失，或出口气体与右侧气体存在温差，\(\Delta S\)
将大于零，输出功也会小于上述最大值。

\newpage

\section{第 6 题}\label{ux7b2c-6-ux9898}

两个半径均为 \(r\) 的同轴圆形线圈位于互相平行的 \(x\)---\(y\)
平面内，圆心分别位于

\[
(0,0,+\lambda r),
\qquad
(0,0,-\lambda r).
\]

两线圈中均通有恒定电流 \(I\)，绕向相同，使它们在轴线上产生的磁场均沿
\(+z\) 方向。

\begin{enumerate}
\def\labelenumi{\arabic{enumi}.}
\tightlist
\item
  求原点的磁感应强度 \(\bm B(0,0,0)\)。
\item
  求原点附近一点 \((x,y,z)\) 的磁场，保留到 \(x/r\)、\(y/r\)、\(z/r\)
  的二阶小量。
\item
  在原点放置一个大小恒为 \(m_0\)、方向只能沿 \(+z\)
  的磁偶极子。若它的中心只能在 \(x\)---\(y\)
  平面内移动，求原点稳定平衡时 \(\lambda\) 应满足的条件。
\item
  若 \(\lambda\) 满足第 3 问条件，改为只允许磁偶极子沿 \(z\)
  轴移动，判断原点平衡是否稳定。
\item
  保持第 3 问的 \(\lambda\)
  条件和磁偶极子的轴向约束，把两恒流线圈改为闭合超导线圈。判断原点能否成为稳定平衡，并说明还需要哪些线圈参数；若忽略两线圈间的互感、每个线圈自感为
  \(L\)，写出稳定条件。
\end{enumerate}

\begin{figure}[H]
\centering
\includegraphics[width=0.73\linewidth,height=\textheight,keepaspectratio,alt={图 6-1　同轴双圆环的坐标、间距和电流绕向。}]{../figures/question-06.pdf}
\caption*{图 6-1　同轴双圆环的坐标、间距和电流绕向。}
\end{figure}

\subsection{一、题型定位}\label{ux4e00ux9898ux578bux5b9aux4f4d-5}

线圈轴线上场强容易求，轴外磁场则用无源区的麦克斯韦方程展开。得到局部场后，再计算磁偶极势能与超导磁通反馈对稳定性的影响。

\subsection{二、知识点}\label{ux4e8cux77e5ux8bc6ux70b9-5}

\begin{itemize}
\item
  圆电流轴线上磁场

  \[B_z(s)=\frac{\mu_0Ir^2}{2(r^2+s^2)^{3/2}}.\]
\item
  对称双线圈轴向磁场及泰勒展开。
\item
  无源区麦克斯韦方程

  \[\nabla\cdot\bm B=0,
  \qquad
  \nabla\times\bm B=0.\]
\item
  轴对称磁场的二阶局部形式。
\item
  固定方向磁偶极势能 \(U=-m_0B_z\)。
\item
  二次势能系数与稳定性的关系。
\item
  超导回路的总磁通守恒 \(L\,\delta I+\delta\Phi_{\rm ext}=0\)。
\item
  互易性：偶极子在线圈处产生的磁通，可由单位线圈电流在偶极子处产生的场求得。
\end{itemize}

\subsection{三、物理图像}\label{ux4e09ux7269ux7406ux56feux50cf-5}

两线圈关于原点镜面对称，原点场沿 \(z\)
轴；所有一次空间导数因对称性为零。原点附近的场强变化首先是二次项。若
\(B_z\) 在横向有局部极大值，磁偶极势能 \(-m_0B_z\)
在横向就是局部极小值；但无源静磁场不能同时在三个方向都形成同号曲率，因此轴向会成为不稳定方向。

超导线圈受扰动后不允许总磁通改变。偶极子向上偏移时，上、下线圈中的外磁通改变量符号相反，于是产生反向的感应电流增量。两环的反馈力可与恒流背景场的轴向失稳力竞争。

\subsection{四、解题入口}\label{ux56dbux89e3ux9898ux5165ux53e3-5}

先把单环轴线场函数记为

\[
b(s)=\frac{\mu_0Ir^2}{2(r^2+s^2)^{3/2}}.
\]

双环轴线场是 \(b(z-\lambda r)+b(z+\lambda r)\)。求出 \(B_0\) 与

\[k=\left.\frac{\mathrm{d}^2B_z}{\mathrm{d}z^2}\right|_{z=0}\]

后，利用轴对称与麦克斯韦方程写出全部二阶展开。稳定性先检查势能二次项；二次项为零时须算更高阶。

\subsection{五、解析}\label{ux4e94ux89e3ux6790-5}

图 6-2 左侧把横向势能 \(U(x,0,0)\) 和轴向势能 \(U(0,0,z)\)
画在同一坐标中，直接显示恒流双环形成鞍点；右侧画出磁偶极子上移时上下超导环的磁通变化与感应电流方向。第（3）---（5）步使用左图，第（6）---（7）步使用右图。

\begin{figure}[H]
\centering
\includegraphics[width=0.9\linewidth,height=\textheight,keepaspectratio,alt={图 6-2　恒流双环的鞍形势能与超导磁通反馈产生的附加回复力。}]{../figures/solution-06.pdf}
\caption*{图 6-2　恒流双环的鞍形势能与超导磁通反馈产生的附加回复力。}
\end{figure}

\textbf{（1）原点磁场} 单个半径为 \(r\) 的圆电流，在距其中心轴向距离
\(s\) 处产生

\[
B_{\rm loop}(s)=\frac{\mu_0Ir^2}{2(r^2+s^2)^{3/2}}.
\]

原点到每个线圈中心的轴向距离大小均为

\[s=\lambda r.\]

每个线圈在原点贡献

\[
\begin{aligned}
B_1
&=\frac{\mu_0Ir^2}
{2[r^2+(\lambda r)^2]^{3/2}}\\
&=\frac{\mu_0Ir^2}
{2r^3(1+\lambda^2)^{3/2}}\\
&=\frac{\mu_0I}{2r(1+\lambda^2)^{3/2}}.
\end{aligned}
\]

两者方向均为 \(+z\)，所以

\[
\boxed{
\bm B(0,0,0)
=\frac{\mu_0I}{r(1+\lambda^2)^{3/2}}\,\hat{\bm z}
}.
\]

记

\[
B_0=\frac{\mu_0I}{r(1+\lambda^2)^{3/2}}.
\]

\textbf{（2）先求轴线上场的二阶导数}

轴线上任意点 \(z\) 的总场为

\[
B_z(0,0,z)=\frac{\mu_0Ir^2}{2}
\left\{
[r^2+(z-\lambda r)^2]^{-3/2}
+[r^2+(z+\lambda r)^2]^{-3/2}
\right\}.
\]

设

\[f(u)=(r^2+u^2)^{-3/2}.\]

第一次求导：

\[
f'(u)=-3u(r^2+u^2)^{-5/2}.
\]

第二次求导时使用乘积法则：

\[
\begin{aligned}
f''(u)
&=-3(r^2+u^2)^{-5/2}
+(-3u)\left[-\frac52(r^2+u^2)^{-7/2}(2u)\right]\\
&=-3(r^2+u^2)^{-5/2}
+15u^2(r^2+u^2)^{-7/2}\\
&=\frac{-3(r^2+u^2)+15u^2}{(r^2+u^2)^{7/2}}\\
&=\frac{-3r^2+12u^2}{(r^2+u^2)^{7/2}}.
\end{aligned}
\]

在原点，两个线圈对应 \(u=\pm\lambda r\)，二阶导数相同。因此

\[
\begin{aligned}
k
&\equiv\left.\frac{\mathrm{d}^2B_z}{\mathrm{d}z^2}\right|_0\\
&=\frac{\mu_0Ir^2}{2}\cdot2
\frac{-3r^2+12\lambda^2r^2}
{[r^2(1+\lambda^2)]^{7/2}}\\
&=\frac{\mu_0Ir^4(12\lambda^2-3)}
{r^7(1+\lambda^2)^{7/2}}\\
&=\boxed{
\frac{\mu_0I}{r^3}
\frac{12\lambda^2-3}{(1+\lambda^2)^{7/2}}
}.
\end{aligned}
\]

轴线场在原点附近为

\[
B_z(0,0,z)=B_0+\frac12kz^2+O(z^4).
\]

\textbf{（3）由麦克斯韦方程恢复三维二阶展开}

为解释图 6-2
左侧为什么一方向上升、另一方向下降，先由原点附近无电流区的麦克斯韦方程恢复三维二阶项。令

\[\rho^2=x^2+y^2.\]

轴对称和镜面对称允许写成

\[
B_z=B_0+az^2+b\rho^2,
\]

\[
B_\rho=c\rho z.
\]

轴线上 \(B_z=B_0+kz^2/2\)，所以

\[a=\frac k2.\]

原点附近没有电流，故 \(\nabla\times\bm B=0\)。在柱坐标中，\(\varphi\)
分量给出

\[
\frac{\partial B_\rho}{\partial z}
-\frac{\partial B_z}{\partial\rho}=0.
\]

分别求导：

\[
\frac{\partial B_\rho}{\partial z}=c\rho,
\qquad
\frac{\partial B_z}{\partial\rho}=2b\rho.
\]

所以

\[c=2b.\]

再使用 \(\nabla\cdot\bm B=0\)：

\[
\frac1\rho\frac{\partial(\rho B_\rho)}{\partial\rho}
+\frac{\partial B_z}{\partial z}=0.
\]

第一项为

\[
\frac1\rho\frac{\partial(c\rho^2z)}{\partial\rho}
=2cz,
\]

第二项为

\[\frac{\partial B_z}{\partial z}=2az.\]

故

\[2c+2a=0,\]

即

\[c=-a=-\frac k2.\]

由 \(c=2b\) 得

\[b=-\frac k4.\]

因此

\[
\boxed{
B_z=B_0+\frac k2z^2-
\frac k4(x^2+y^2)
}.
\]

柱坐标横向分量为

\[
B_\rho=-\frac k2\rho z.
\]

换成直角坐标：

\[
\boxed{B_x=-\frac k2xz},
\]

\[
\boxed{B_y=-\frac k2yz}.
\]

所以二阶局部磁场完整写为

\[
\boxed{
\bm B=
-\frac k2xz\,\hat{\bm x}
-\frac k2yz\,\hat{\bm y}
+\left[
B_0+\frac k2z^2-\frac k4(x^2+y^2)
\right]\hat{\bm z}
}.
\]

\textbf{（4）磁偶极子在 \(x\)---\(y\) 平面内的稳定条件}

按图 6-2 左侧蓝色横向势阱，磁偶极矩固定为

\[\bm m=m_0\hat{\bm z}.\]

势能为

\[U=-\bm m\cdot\bm B=-m_0B_z.\]

在 \(z=0\) 平面内：

\[
B_z(x,y,0)=B_0-\frac k4(x^2+y^2).
\]

因此

\[
\begin{aligned}
U(x,y,0)
&=-m_0B_0+\frac{m_0k}{4}(x^2+y^2).
\end{aligned}
\]

原点成为势能极小值，要求二次项系数为正：

\[m_0k>0.\]

题设 \(m_0>0\)，所以要求

\[k>0.\]

由 \(k\) 的表达式，分母始终为正，只需

\[12\lambda^2-3>0.\]

逐步化简：

\[12\lambda^2>3,\]

\[\lambda^2>\frac14.\]

由于 \(\lambda\) 表示正距离参数，故

\[
\boxed{\lambda>\frac12}.
\]

\(\lambda=1/2\) 时二次项为零，二阶展开本身不能判断。继续计算得到

\[\frac{B_z(0,0,z)}{B_0}=1-\frac{144}{125}\left(\frac zr\right)^4+O(z^6/r^6),\]

\[\frac{B_z(x,y,0)}{B_0}=1-\frac{54}{125}\left(\frac{\sqrt{x^2+y^2}}r\right)^4+O((x^2+y^2)^3/r^6).\]

两种单独受限方向的势能四次项均为正，原点均为稳定平衡。因此第3问的完整条件为
\(\boxed{\lambda\ge1/2}\)；\(\lambda>1/2\)
是具有非零线性回复力的条件。二维或三维混合方向仍不能据此判为稳定。

\textbf{（5）只允许沿 \(z\) 轴运动时的稳定性}

图中 \(\lambda>1/2\) 的红色轴向曲线在原点为势垒顶。临界 \(\lambda=1/2\)
须使用上一段四阶式。在轴线上

\[
B_z(0,0,z)=B_0+\frac k2z^2.
\]

势能为

\[
U(0,0,z)=-m_0B_0-\frac{m_0k}{2}z^2.
\]

当 \(\lambda>1/2\) 时，\(k>0\)，此时 \(z^2\)
项系数为负，原点是势能极大值。受力也可直接看出：

\[
F_z=-\frac{\mathrm{d}U}{\mathrm{d}z}=m_0kz.
\]

若 \(z>0\)，则 \(F_z>0\)，力把偶极子继续推向正方向；若 \(z<0\)，则
\(F_z<0\)，仍把它推离原点。因此

\[
\boxed{\lambda>1/2\text{ 时轴向不稳定；}\lambda=1/2\text{ 时四阶稳定。}}
\]

\textbf{（6）超导线圈的磁通反馈}

转看图 6-2 右侧：磁偶极子向上扰动 \(z\)
后，上、下两环磁通变化符号相反，感应电流也反向。把恒流线圈改为闭合超导线圈后，电流不再固定。理想超导回路保持总磁通不变：

\[L\,\delta I+\delta\Phi_m=0.\]

先求偶极子移动引起的外磁通变化。令单位电流圆环在轴向距离 \(s\)
处产生的场为

\[
b(s)=\frac{\mu_0r^2}{2(r^2+s^2)^{3/2}}.
\]

由电磁互易性，磁矩 \(m_0\) 对该圆环贡献的磁通可写成

\[\Phi_m(s)=m_0b(s).\]

求导：

\[
b'(s)=-\frac{3\mu_0r^2s}
{2(r^2+s^2)^{5/2}}.
\]

偶极子从原点向 \(+z\) 偏移一个小量
\(z\)。对上方线圈，偶极子相对线圈的轴向坐标为 \(s=z-\lambda r\)。在
\(z=0\) 时

\[
b'(-\lambda r)
=\frac{3\mu_0\lambda}
{2r^2(1+\lambda^2)^{5/2}}.
\]

定义

\[
A=m_0b'(-\lambda r)
=\frac{3\mu_0m_0\lambda}
{2r^2(1+\lambda^2)^{5/2}}.
\]

于是上、下线圈的外磁通变化分别为

\[\delta\Phi_U=+Az,\]

\[\delta\Phi_L=-Az.\]

忽略两线圈互感时，磁通守恒给出

\[
\delta I_U=-\frac{A}{L}z,
\qquad
\delta I_L=+\frac{A}{L}z.
\]

两线圈的电流增量方向相反。

\textbf{（7）计算感应电流产生的轴向回复力}

偶极子与一个线圈的相互作用力为

\[F_z=m_0I\,b'(s).\]

只保留由感应电流产生、且与 \(z\) 成正比的项。上方线圈贡献

\[
\delta F_U
=m_0\delta I_U b'(-\lambda r).
\]

因 \(A=m_0b'(-\lambda r)\)，有

\[
\delta F_U
=\left(-\frac A Lz\right)A
=-\frac{A^2}{L}z.
\]

对下方线圈，\(b'(+\lambda r)=-A/m_0\)，而 \(\delta I_L=+Az/L\)，所以

\[
\delta F_L
=m_0\left(\frac A Lz\right)
\left(-\frac A{m_0}\right)
=-\frac{A^2}{L}z.
\]

总反馈力为

\[
\boxed{
F_{\rm ind}=-\frac{2A^2}{L}z
}.
\]

恒流背景磁场原有的失稳力为

\[
F_{\rm bg}=m_0kz.
\]

总线性力为

\[
F_z=\left(m_0k-\frac{2A^2}{L}\right)z.
\]

稳定要求 \(z>0\) 时 \(F_z<0\)，即

\[
\frac{2A^2}{L}>m_0k.
\]

代入 \(A\) 与 \(k\)：

\[
\frac{2}{L}
\left[
\frac{3\mu_0m_0\lambda}
{2r^2(1+\lambda^2)^{5/2}}
\right]^2
>
m_0\frac{\mu_0I}{r^3}
\frac{12\lambda^2-3}{(1+\lambda^2)^{7/2}}.
\]

左边平方并化简：

\[
\frac{9\mu_0^2m_0^2\lambda^2}
{2Lr^4(1+\lambda^2)^5}
>
\frac{\mu_0Im_0(12\lambda^2-3)}
{r^3(1+\lambda^2)^{7/2}}.
\]

两边除以正量 \(\mu_0m_0/[r^3(1+\lambda^2)^{7/2}]\)：

\[
\boxed{
\frac{9\mu_0m_0\lambda^2}
{2Lr(1+\lambda^2)^{3/2}}
>I(12\lambda^2-3)
}.
\]

等价地，

\[
\boxed{
LI<
\frac{9\mu_0m_0\lambda^2}
{2r(1+\lambda^2)^{3/2}(12\lambda^2-3)}
},
\qquad \lambda>\frac12.
\]

所以超导反馈有可能把轴向失稳改为稳定，但不能仅凭 \(\lambda\)
判定；还必须知道自感 \(L\)、背景电流 \(I\)、偶极矩 \(m_0\) 和几何尺度
\(r\)。若两环同向绕制且互感为 \(M\)，反对称电流模的电感为
\(L-M>0\)，总回复项改为 \(2A^2/(L-M)\)。在 \(\lambda=1/2\) 时
\(k=0\)，有限正电感的超导反馈给出正线性刚度，轴向稳定；上面的除以
\(12\lambda^2-3\) 的不等式只适用于 \(\lambda>1/2\)。

\textbf{题面坐标的几何校正}

题面文字把圆心坐标写成 \((\pm\lambda r,0,0)\)，同时又声明圆环位于
\(x\)---\(y\) 平面。若照此理解，两圆环圆心沿 \(x\)
方向错开，便不再同轴，也与配图和后续轴向稳定问题矛盾。根据图形、轴对称性及后续运动条件，圆心应位于
\((0,0,\pm\lambda r)\)，下文按这一自洽几何计算。

\textbf{超导反馈中因子 \(2\) 的来源}

上、下两个超导环各自产生 \(-A^2z/L\) 的回复力，因此总反馈为
\(-2A^2z/L\)。若只计算一只线圈或在能量法中没有正确处理磁通约束，就会漏掉因子
\(2\)；逐一计算两线圈的力可直接排除这一错误。

\textbf{互感的影响}

扰动引起 \(\delta I_U=-\delta I_L\)，属于反对称电流模。若两线圈互感为
\(M\)，磁链方程给出的有效电感通常为 \(L-M\)，而不是单环
\(L\)。具体正负号取决于线圈绕向和互感符号约定，所以题目未给 \(M\)
时不能声称得到完全一般的超导稳定条件。

\newpage

\section{第 7 题}\label{ux7b2c-7-ux9898}

一束近似平行光沿 \(+y\) 方向入射到柱形透镜。透镜的一面为平面，位于
\(x\)---\(z\) 平面；另一面在 \(x\)---\(y\) 截面内为半径

\[r=5.00\ \mathrm{cm}\]

的半圆弧。透镜折射率为
\(n=1.56\)，外界为空气。平行光垂直入射平面，因此只在半圆弧面发生折射。弧面上方铺有与光路近似平行的粗糙白纸，可观察实际光路。透镜圆心取为原点。

\begin{enumerate}
\def\labelenumi{\arabic{enumi}.}
\tightlist
\item
  近轴光线会聚于焦点 \(F\)，求焦点到原点的距离 \(h_0\)。
\item
  忽略多次反射。入射平行光距 \(y\) 轴的最大横向距离 \(\xi_{\max}\)
  为多少时，弧面处仍有折射光射出？
\item
  在过焦点 \(F\) 且垂直于 \(y\) 轴的截面上，实际光束沿 \(x\)
  方向的几何宽度最大为多少？
\item
  求 \(y=10.0\ \mathrm{cm}\) 平面上被照亮区域沿 \(x\) 方向的宽度。
\end{enumerate}

\begin{figure}[H]
\centering
\includegraphics[width=0.53\linewidth,height=\textheight,keepaspectratio,alt={图 7-1　平行光垂直入射半圆柱透镜的截面光路。}]{../figures/question-07.pdf}
\caption*{图 7-1　平行光垂直入射半圆柱透镜的截面光路。}
\end{figure}

\subsection{一、题型定位}\label{ux4e00ux9898ux578bux5b9aux4f4d-6}

近轴光线先给出理想焦距，有限口径的边缘光线会因曲面折射和全反射改变范围。几问使用的光线模型不同，不能把焦点公式直接用于边缘束宽。

\subsection{二、知识点}\label{ux4e8cux77e5ux8bc6ux70b9-6}

\begin{itemize}
\item
  平面垂直入射不改变传播方向。
\item
  近轴折射关系与单球面焦距。
\item
  斯涅耳定律

  \[n\sin\theta_1=\sin\theta_2.\]
\item
  从高折射率介质射向空气的临界角

  \[\sin\theta_c=\frac1n.\]
\item
  圆弧点坐标 \((r\sin\theta,r\cos\theta)\)。
\item
  临界折射光沿圆弧切线传播。
\item
  有限口径球差及边界光线法。
\end{itemize}

\subsection{三、物理图像}\label{ux4e09ux7269ux7406ux56feux50cf-6}

入射光在平面处沿法线进入，不偏折。到达圆弧面后，法线就是圆心到出射点的半径。近轴光线入射角很小，折射后在轴上近似会聚于同一点；远轴光线偏折更强，不再通过同一焦点。达到临界入射角的最外侧光线沿弧面切线射出，它构成传播区域的一条重要边界。

\subsection{四、解题入口}\label{ux56dbux89e3ux9898ux5165ux53e3-6}

令弧面出射点相对 \(y\) 轴的横坐标为 \(\xi\)，法线相对 \(+y\) 轴的夹角为
\(\theta_1\)。几何关系为

\[
\sin\theta_1=\frac\xi r.
\]

近轴时线性化斯涅耳定律；临界时令折射角
\(\theta_2=90^\circ\)。在给定水平截面上，最外边界由临界切线与该截面的交点给出。

\subsection{五、解析}\label{ux4e94ux89e3ux6790-6}

图 7-2 同时标出近轴焦点
\(F(0,h_0)\)、弧面临界出射点、与弧面相切的边界光线以及两个观察平面。第（1）步确定橙色近轴焦线，第（2）---（3）步确定红色边界光线，第（4）---（5）步读取它与指定平面的交点宽度。

\begin{figure}[H]
\centering
\includegraphics[width=0.566\linewidth,height=\textheight,keepaspectratio,alt={图 7-2　近轴焦点、临界入射点和切向边界光线。}]{../figures/solution-07.pdf}
\caption*{图 7-2　近轴焦点、临界入射点和切向边界光线。}
\end{figure}

\textbf{（1）近轴焦点位置} 先沿图 7-2 的橙色近轴光线求轴上交点
\(F(0,h_0)\)。

取一条距轴很近的光线，其弧面入射点横坐标为 \(\xi\)，且
\(\xi\ll r\)。入射点的纵坐标为

\[
y_s=\sqrt{r^2-\xi^2}\simeq r.
\]

法线与 \(+y\) 轴的夹角为 \(\theta_1\)，小角下

\[
\sin\theta_1=\frac\xi r
\quad\Longrightarrow\quad
\theta_1\simeq\frac\xi r.
\]

在弧面处，光从折射率 \(n\) 的介质射向空气。斯涅耳定律为

\[n\sin\theta_1=\sin\theta_2.\]

近轴小角下

\[n\theta_1\simeq\theta_2.\]

出射光相对 \(y\) 轴的偏折角为

\[
\delta=\theta_2-\theta_1
\simeq(n-1)\theta_1.
\]

代入 \(\theta_1\simeq\xi/r\)：

\[
\delta\simeq(n-1)\frac\xi r.
\]

从弧面顶点附近到焦点的距离记为 \(f\)。小角几何给出

\[
\tan\delta\simeq\delta\simeq\frac\xi f.
\]

所以

\[
\frac\xi f=(n-1)\frac\xi r.
\]

对非轴线上光线约去 \(\xi\)：

\[
\boxed{f=\frac r{n-1}}.
\]

圆心到弧面顶点的距离为 \(r\)，所以焦点到原点的距离为

\[
\begin{aligned}
h_0
&=r+f\\
&=r+\frac r{n-1}\\
&=\frac{nr}{n-1}.
\end{aligned}
\]

代入 \(r=5.00\ \mathrm{cm}\)、\(n=1.56\)：

\[
\begin{aligned}
h_0
&=\frac{1.56\times5.00}{1.56-1}\ \mathrm{cm}\\
&=\frac{7.80}{0.56}\ \mathrm{cm}\\
&\simeq13.93\ \mathrm{cm}.
\end{aligned}
\]

因此

\[\boxed{h_0\simeq13.93\ \mathrm{cm}}.\]

\textbf{（2）仍能从弧面折射出的最大入射高度} 图 7-2
红色边界光线从弧面临界点出射；该点的入射角恰好对应折射角 \(90^\circ\)。

光从高折射率 \(n\) 射向空气。临界状态满足

\[
n\sin\theta_c=1.
\]

所以

\[
\boxed{\sin\theta_c=\frac1n}.
\]

弧面上一点横坐标为

\[\xi=r\sin\theta_1.\]

最大可出射光对应 \(\theta_1=\theta_c\)：

\[
\xi_{\max}=r\sin\theta_c=\frac rn.
\]

数值为

\[
\begin{aligned}
\xi_{\max}
&=\frac{5.00}{1.56}\ \mathrm{cm}\\
&\simeq3.205\ \mathrm{cm}.
\end{aligned}
\]

故

\[\boxed{\xi_{\max}\simeq3.21\ \mathrm{cm}}.\]

超过这一横向距离的光线在弧面发生全反射；题目要求忽略多次反射，因此不计其后续传播。

\textbf{（3）写出临界边界光线方程} 在图 7-2
中，临界出射光线与圆弧相切，因此可用临界点坐标和切线斜率直接写方程。

临界点坐标为

\[
x_c=r\sin\theta_c,
\qquad
y_c=r\cos\theta_c.
\]

临界折射角为
\(90^\circ\)，所以出射光沿圆弧切线。对右侧临界点，向上的切线指向左上方。切线相对
\(+y\) 轴的夹角为

\[90^\circ-\theta_c.\]

因此其斜率关系为

\[
\frac{\mathrm{d}x}{\mathrm{d}y}=-\cot\theta_c.
\]

经过 \((x_c,y_c)\) 的边界光线为

\[
\boxed{
x(y)=r\sin\theta_c
-(y-r\cos\theta_c)\cot\theta_c
}.
\]

左侧存在关于 \(y\) 轴对称的边界光线。

\textbf{（4）焦点平面处的最大宽度}

在 \(y=h_0\) 平面，右侧临界光线已经穿过轴线并到达负 \(x\)
一侧，其横坐标为

\[
x_F=r\sin\theta_c
-(h_0-r\cos\theta_c)\cot\theta_c.
\]

此数为负。半宽取其绝对值：

\[
a_F=(h_0-r\cos\theta_c)\cot\theta_c-r\sin\theta_c.
\]

总宽度为

\[
\boxed{
D_F=2\left[
(h_0-r\cos\theta_c)\cot\theta_c
-r\sin\theta_c
\right]
}.
\]

数值计算。先有

\[
\sin\theta_c=\frac1{1.56}\simeq0.64103,
\]

\[
\cos\theta_c
=\sqrt{1-\frac1{1.56^2}}
\simeq0.76752,
\]

\[
\cot\theta_c
=\frac{\cos\theta_c}{\sin\theta_c}
\simeq1.1973.
\]

于是

\[
\begin{aligned}
a_F
&=(13.9286-5.00\times0.76752)\times1.1973
-5.00\times0.64103\\
&\simeq8.877\ \mathrm{cm}.
\end{aligned}
\]

所以

\[
\boxed{D_F\simeq17.75\ \mathrm{cm}}.
\]

该宽度远大于透镜直径并不矛盾：临界边界光在越过轴线后继续向外发散，焦点平面只对近轴光线是会聚平面。

\textbf{（5）\(y=10.0\ \mathrm{cm}\) 平面上的照亮宽度}

把 \(y=10.0\ \mathrm{cm}\) 代入边界光线。右侧临界光线的横坐标为

\[
x_{10}=r\sin\theta_c
-(10.0-r\cos\theta_c)\cot\theta_c.
\]

代入数值：

\[
\begin{aligned}
x_{10}
&=5.00\times0.64103
-(10.0-5.00\times0.76752)\times1.1973\\
&\simeq-4.173\ \mathrm{cm}.
\end{aligned}
\]

左右对称，照亮区域宽度为

\[
D_{10}=2|x_{10}|.
\]

因此

\[\boxed{D_{10}\simeq8.35\ \mathrm{cm}}.\]

\textbf{一般出射光线的核验式}

对任意 \(0\le\xi\le r/n\)，定义

\[
\theta_1=\arcsin\frac\xi r,
\qquad
\theta_2=\arcsin\left(n\frac\xi r\right).
\]

出射光相对轴线向内偏折

\[\delta=\theta_2-\theta_1.\]

其光线方程为

\[
x(y;\xi)=\xi-
\left(y-\sqrt{r^2-\xi^2}\right)\tan\delta.
\]

令 \(\xi\to0\) 并保留一阶，交轴位置趋于 \(h_0\)；令
\(\xi\to r/n\)，\(\theta_2\to\pi/2\)，上式退化为临界切线方程。这两个极限共同验证了近轴与临界两部分推导。

\textbf{数值辨认与回代}

第三问的数值应为 \(17.76\ \mathrm{cm}\)，不应误读成
\(7.76\ \mathrm{cm}\)。直接代入公式得到
\(17.75\ \mathrm{cm}\)，也验证了其数量级和首位数字。

\newpage

\section{第 8 题}\label{ux7b2c-8-ux9898}

一个半径为 \(R\)
的半球形光滑碗固定在水平地面上，对称轴竖直。一个小球与碗的内表面接触，初始位置比碗底高

\[\frac R3.\]

\begin{enumerate}
\def\labelenumi{\arabic{enumi}.}
\tightlist
\item
  给小球一个与碗面相切、方向可任意选择的初速度。求小球能够越过碗沿所需的最小初速度，并说明达到该最小值的发射方向有几种。
\item
  若初速度限定在水平方向且与碗面相切，求小球能够越过碗沿所需的最小速度。
\end{enumerate}

\begin{figure}[H]
\centering
\includegraphics[width=0.892\linewidth,height=\textheight,keepaspectratio,alt={图 8-1　半球碗内小球的初始高度、任意切向速度与水平切向速度。}]{../figures/question-08.pdf}
\caption*{图 8-1　半球碗内小球的初始高度、任意切向速度与水平切向速度。}
\end{figure}

\subsection{一、题型定位}\label{ux4e00ux9898ux578bux5b9aux4f4d-7}

初始切向方向可以选择时，应分别考察子午与方位运动对能量的贡献。若速度沿水平切线，方位角动量则由轴对称性守恒。

\subsection{二、知识点}\label{ux4e8cux77e5ux8bc6ux70b9-7}

\begin{itemize}
\item
  光滑固定约束力不做功，机械能守恒。
\item
  轴对称系统绕竖直轴无外力矩，\(L_z\) 守恒。
\item
  球面上某高度处的水平圆半径

  \[\rho(h)=\sqrt{R^2-(R-h)^2}.\]
\item
  子午方向与方位方向的速度分解。
\item
  有效势能

  \[U_{\rm eff}(h)=mgh+\frac{L_z^2}{2m\rho(h)^2}.\]
\item
  阈值运动在碗沿处的子午速度为零。
\item
  约束力非负条件的检查。
\end{itemize}

\subsection{三、物理图像}\label{ux4e09ux7269ux7406ux56feux50cf-7}

碗面上的切平面有两个互相垂直的基本方向：沿经线升降的子午方向，以及沿等高圆绕轴转动的方位方向。重力只改变高度，方位速度却因
\(L_z\)
守恒不能随意消失。为了用最少初动能越过碗沿，应尽量不把能量留在不可消去的方位运动中。

\subsection{四、解题入口}\label{ux56dbux89e3ux9898ux5165ux53e3-7}

第一问先由能量给出任何路径都必须满足的下界，再寻找能实现该下界的方向。第二问计算初始位置的水平圆半径
\(\rho_0\)，由

\[m\rho_0v_0=mRv_{\varphi,\rm rim}\]

求碗沿仍保留的方位速度，然后代入能量守恒。

\subsection{五、解析}\label{ux4e94ux89e3ux6790-7}

图 8-2 左侧画出沿子午线向上的两条最省能路径，它们都满足轴角动量
\(L_z=0\)；右侧画出水平切向发射时包含重力势能和方位离心项的有效势。第（1）---（2）步使用左图，第（3）---（6）步使用右图。

\begin{figure}[H]
\centering
\includegraphics[width=0.9\linewidth,height=\textheight,keepaspectratio,alt={图 8-2　零轴角动量的两条子午路径与水平发射时的有效势能。}]{../figures/solution-08.pdf}
\caption*{图 8-2　零轴角动量的两条子午路径与水平发射时的有效势能。}
\end{figure}

\textbf{（1）任意切向发射时的能量下界} 按图 8-2
左侧，无论选择近侧还是对侧子午路径，到达碗沿都至少要克服同一重力势能差。以碗底为重力势能零点，小球初始高度为

\[h_0=\frac R3.\]

碗沿高度为

\[h_{\rm rim}=R.\]

小球至少需要克服的重力势能增量为

\[
\Delta U=mg(R-R/3)=\frac23mgR.
\]

若初速度为 \(v_0\)，初动能为

\[K_0=\frac12mv_0^2.\]

无论选择何种路径，都必须有

\[
\frac12mv_0^2\ge\frac23mgR.
\]

约去 \(m\)，两边乘 \(2\)：

\[
v_0^2\ge\frac43gR.
\]

所以能量给出的绝对下界为

\[
v_0\ge2\sqrt{\frac{gR}{3}}.
\]

\textbf{（2）证明下界可以达到，并判断方向数}

图 8-2
左侧蓝、绿两条路径都没有方位分量。若初速度含有方位分量，则绕竖直轴的角动量

\[L_z=m\rho v_\varphi\]

不为零。到达碗沿时 \(\rho=R\)，仍必须保留方位动能

\[K_{\varphi,\rm rim}=\frac{L_z^2}{2mR^2}>0.\]

这样初动能不仅要补足势能增量，还要留下非零动能，因而不能达到上述能量下界。

要达到最小值，必须令

\[L_z=0.\]

在初始点的切平面内，\(L_z=0\)
等价于速度没有方位分量，即速度沿通过竖直轴的子午线方向。沿同一条子午线有两个相反方向：

\begin{itemize}
\tightlist
\item
  沿近侧碗面直接向上，到达近侧碗沿；
\item
  先沿碗面向下经过碗底，再从对侧上升到对侧碗沿。
\end{itemize}

两条路径都能在阈值下到达碗沿，且此时速度为零。若``越过''要求实际跨出碗沿，子午速度还要严格为正：阈值等号到达后会回落，所求是逃逸速度的下确界。两个最优发射方向仍是这两个相反的子午方向。

所以

\[
\boxed{v_{\min}=2\sqrt{\frac{gR}{3}}},
\]

达到最小值的方向共有

\[\boxed{2\ \text{种}}.\]

\textbf{（3）水平切向发射时的初始水平圆半径}

转看图 8-2
右侧的有效势。水平切向发射时轴角动量不再为零。半球球心位于碗底上方
\(R\)，初始点相对球心的竖直坐标为

\[
z_0=h_0-R
=\frac R3-R
=-\frac{2R}{3}.
\]

球面方程给出

\[
\rho_0^2+z_0^2=R^2.
\]

所以

\[
\begin{aligned}
\rho_0
&=\sqrt{R^2-z_0^2}\\
&=\sqrt{R^2-\frac{4R^2}{9}}\\
&=R\sqrt{\frac59}\\
&=\frac{\sqrt5}{3}R.
\end{aligned}
\]

即

\[\boxed{\rho_0=\frac{\sqrt5}{3}R}.\]

\textbf{（4）用轴角动量守恒求碗沿方位速度}

题设初速度水平且与碗面相切。在球面上，这正是沿等高圆的方位方向，所以初始速度全部为
\(v_{\varphi0}=v_0\)。

绕竖直轴角动量守恒：

\[
m\rho_0v_0=mRv_{\varphi,\rm rim}.
\]

约去 \(m\)，代入 \(\rho_0=\sqrt5R/3\)：

\[
\frac{\sqrt5}{3}Rv_0
=Rv_{\varphi,\rm rim}.
\]

约去 \(R\)：

\[
\boxed{
v_{\varphi,\rm rim}=\frac{\sqrt5}{3}v_0
}.
\]

阈值状态下，小球刚好到达碗沿，子午方向速度为零，但上述方位速度不能为零。此刻子午加速度仍朝碗底，接触会继续保持，随后回落；实际越出要求严格超过阈值。

\textbf{（5）机械能守恒求水平发射阈值}

初态机械能为

\[
E_i=\frac12mv_0^2+mg\frac R3.
\]

临界到达碗沿时的机械能为

\[
E_f=\frac12m v_{\varphi,\rm rim}^2+mgR.
\]

令 \(E_i=E_f\)：

\[
\frac12mv_0^2+\frac13mgR
=\frac12m\left(\frac{\sqrt5}{3}v_0\right)^2+mgR.
\]

平方方位速度：

\[
\left(\frac{\sqrt5}{3}v_0\right)^2
=\frac59v_0^2.
\]

代回并约去 \(m\)：

\[
\frac12v_0^2+\frac13gR
=\frac{5}{18}v_0^2+gR.
\]

把速度项移到左边、势能项移到右边：

\[
\left(\frac12-\frac5{18}\right)v_0^2
=\left(1-\frac13\right)gR.
\]

分别通分：

\[
\frac{9-5}{18}v_0^2=\frac23gR.
\]

所以

\[
\frac29v_0^2=\frac23gR.
\]

两边乘 \(9/2\)：

\[v_0^2=3gR.\]

因此

\[
\boxed{v_0>\sqrt{3gR}}.
\]

\textbf{（6）接触条件检查}

小球在碗内表面运动时，碗面对小球的支持力指向球心。设小球高度为
\(h<R\)、速率为 \(v\)，沿球心径向列式可得

\[
N=\frac{mv^2}{R}+mg\frac{R-h}{R}.
\]

在碗内 \(0\le h<R\)
时两项均非负，所以运动途中不会因所需支持力为负而提前脱离。到达碗沿
\(h=R\) 时

\[N=\frac{mv^2}{R}\ge0.\]

第一问临界路径在碗沿 \(v\to0\)，对应 \(N\to0\)；第二问仍有方位速度，所以
\(N>0\)。两种阈值均与接触约束相容，但到达边缘与跨出边缘应按上述严格性区分。

\textbf{有效势能的统一写法}

对给定的 \(L_z\)，速度可分为子午分量与方位分量：

\[
K=\frac12mv_{\rm mer}^2
+\frac{L_z^2}{2m\rho(h)^2}.
\]

因此高度方向面对的有效势能为

\[
U_{\rm eff}(h)=mgh+
\frac{L_z^2}{2m\rho(h)^2},
\]

其中

\[
\rho(h)^2=R^2-(R-h)^2=2Rh-h^2.
\]

第一问通过选方向令 \(L_z=0\)，有效势能只剩 \(mgh\)；第二问的
\(L_z=m\rho_0v_0\)
由初速同时决定，把碗沿有效势能与初能量相等即可得到同一结果。

\textbf{为什么只有两个最优方向}

在给定初始点的二维切平面中，\(L_z=0\)
是一条通过原点的直线方向，而速度方向是有向的，所以该直线给出两个相反的速度方向。它们分别对应同一子午大圆的两个行进方向，不是连续无穷多个方向。

\end{document}
